\documentclass[a4paper,fleqn]{cas-sc}

\usepackage[authoryear]{natbib}
\usepackage{amsmath,amssymb}
\usepackage{bm}
\usepackage{array}
\usepackage{booktabs}
\usepackage[table]{xcolor}
\usepackage{multirow}
\usepackage{makecell}
\usepackage[caption=false,font=normalsize,labelfont=sf,textfont=sf]{subfig}
\usepackage{textcomp}
\usepackage{url}
\usepackage{graphicx}
\usepackage{tabularx}
\usepackage{algorithmic}
% Improve line breaking to avoid overfull boxes (e.g. long e-mail strings)
\sloppy
\emergencystretch=3em

\begin{document}
\let\WriteBookmarks\relax
%\def\floatpagepagefraction{1}
%\def\textpagefraction{.001}


\shorttitle{QAD-MultiGuard: Edge--Cloud Multimodal Fraud Detection}
\shortauthors{D. Wang et al.}

\title[mode=title]{QAD-MultiGuard: An Edge--Cloud Framework for Multimodal Fraud Detection and Risk Assessment}

\author[1]{Dajin Wang}
\ead{dj_wang@lzre.edu.cn}
\affiliation[1]{organization={School of Information Engineering, Lanzhou Resources \& Environment Vocational Technology University},
            city={Lanzhou},
            postcode={730021},
            country={China}}

\author[2,3]{Jianming Bai}
\ead{baijm@lzu.edu.cn}
\affiliation[2]{organization={School of Management, Lanzhou University},
            city={Lanzhou},
            postcode={730000},
            country={China}}
\affiliation[3]{organization={Center for Data Science, Lanzhou University},
            city={Lanzhou},
            postcode={730000},
            country={China}}

\author[4]{Wu Zhang}
\cormark[1]
\ead{zhangwu@lzre.edu.cn}
\affiliation[4]{organization={School of Management, Lanzhou Resources \& Environment Vocational Technology University},
            city={Lanzhou},
            postcode={730021},
            country={China}}

\author[4]{Wenbin Guo}
\ead{guowenbin@lzre.edu.cn}

\author[4]{Qian Zhao}
\ead{zhaoqian@lzre.edu.cn}

\author[5]{Liangdong Yang}
\ead{yld610158@lzre.edu.cn}
\affiliation[5]{organization={Department of Basic Sciences, Lanzhou Resources \& Environment Vocational Technology University},
            city={Lanzhou},
            postcode={730021},
            country={China}}

\cortext[1]{Corresponding author.}

\nonumnote{This work was supported by the Gansu Provincial Department of Education, China, through the Higher Education Teachers' Innovation Fund Project [grant number 2023A-224]. Jianming Bai's work was supported by the National Natural Science Foundation of China [grant number 72273059] and the National Key Research and Development Program of China [grant number 2024YFB3311600].}

\begin{abstract}
Telecommunication fraud forensics requires real-time reasoning over heterogeneous voice and text signals on commodity mobile hardware under strict privacy regulation, yet existing systems remain predominantly single-modal and degrade markedly in realistic deployment conditions. Real-time fraud detection on resource-constrained mobile devices is governed by three coupled constraints: \textit{privacy} (raw audio and text must remain on-device), \textit{fidelity} (low-bit quantisation must not sacrifice accuracy), and \textit{responsiveness} (per-window inference must stay below 500~ms). This paper presents \textbf{QAD-MultiGuard}, an edge--cloud collaborative framework for audio--text fraud detection that addresses these constraints through four coordinated components: (i) a pure-KL quantisation-aware distillation (QAD) scheme that transfers knowledge from a full-precision BF16 teacher to heterogeneous low-bit students without a separate quantisation-aware training stage; (ii) an Output-Variance Freeze (OV-Freeze) regulariser that stabilises quantisation-sensitive projection layers; (iii) a reconstruction-resistant 128-dimensional acoustic embedding that resists speech-content reconstruction while keeping raw audio on-device; and (iv) a domain-adapted speculative decoding module that accelerates the asynchronous cloud-side chain-of-thought review. On the TeleAntiFraud-28k benchmark, the on-device \texttt{Q4\_K\_M} student attains an $F_1$-score of 0.917 (98.5\% of the BF16 teacher) at a median per-window latency of 268~ms on a Qualcomm Snapdragon 8 Gen 3 platform, while the cloud-side \texttt{NVFP4} track reaches 0.923 (99.1\%) with asynchronous chain-of-thought review, jointly delivering a favourable trade-off among accuracy, privacy, and operational efficiency. The decision-fusion framework is four-modal (text, acoustic, URL, and metadata); the primary TAF-28k evaluation is exercised on the text and acoustic scores available in the corpus, with the URL and metadata fusion weights carried as deployment parameters. Beyond in-domain accuracy, the framework retains robustness under adversarial rewrites (\textit{AdvFraud-3k}) and cross-domain transfer to a text-only Chinese fraud benchmark (\textit{ChiFraud}). These results demonstrate that low-bit quantisation-aware distillation combined with a reconstruction-resistant acoustic embedding can deliver near-real-time multimodal fraud detection on commodity mobile hardware at a small accuracy cost, establishing a feasibility baseline for privacy-oriented on-device AI under strict latency budgets. The reported cloud-side \texttt{NVFP4} figures are produced under a numerical-behaviour-emulation protocol on H100 hardware rather than native Blackwell execution, and the primary tables are drawn from a formal H100 run whose public-repository reproduction is in progress; generalisation to unconstrained real-world deployment remains to be validated through field studies.
\end{abstract}

\begin{highlights}
\item Pure-KL distillation recovers 98.5\% of BF16 accuracy at 4-bit precision.
\item OV-Freeze stabilises training under aggressive quantisation.
\item Reconstruction-resistant embeddings hinder speech-content reconstruction.
\item Domain-adapted speculative decoding yields a $3.32\times$ decoding speedup on Snapdragon 8 Gen 3 for the cloud-side review draft model.
\item On-device student attains $F_1=0.917$ at 268~ms median latency on the TAF-28k benchmark.
\end{highlights}

\begin{keywords}  
Intelligent expert system
\sep multimodal fraud detection
\sep edge--cloud systems 
\sep quantisation-aware distillation 
\sep on-device inference 
\sep reconstruction-resistant representation
\end{keywords}

\maketitle

\section{Introduction}

Designing intelligent frameworks that operate under concurrent real-world constraints---restricted compute, strict privacy regulations, and hard latency budgets---remains an open challenge in applied artificial intelligence. Telecommunication fraud forensics exemplifies this challenge, requiring real-time reasoning over heterogeneous signals---predominantly voice and text, with URL and metadata as auxiliary channels---on commodity mobile hardware under constraints informed by data-protection regulations such as China's Personal Information Protection Law (PIPL), which imposes stringent restrictions on the transmission of raw biometric data.

Recent studies have shown that fraud cases in the TeleAntiFraud-28k corpus exhibit strong multimodal coupling, with each instance involving multiple heterogeneous signals, including voice calls, SMS phishing, malicious URLs, and social-engineering scripts~\citep{1}. The decision-fusion framework is four-modal (text, acoustic, URL, and metadata); the present TAF-28k evaluation is exercised on the text and acoustic scores, which are the only modality scores provided by the corpus, with the URL and metadata fusion weights carried as deployment parameters. However, existing systems remain predominantly single-modal or rely on conventional late-fusion mechanisms. In particular, detection pipelines that rely on automatic speech recognition (ASR) transcripts suffer marked performance degradation in realistic deployment conditions.

For example, SAFE-QAQ, a representative deployed baseline, reports that ASR transcript-based fraud detection suffers marked performance degradation under dialect variation, acoustic noise, and intentional voice manipulation~\citep{2}. These findings highlight the inherent limitations of relying on intermediate textual representations and motivate direct multimodal fusion under joint latency and privacy constraints.

Conventional rule-based and machine learning approaches cannot simultaneously satisfy the concurrent requirements of real-time operation, privacy preservation, and deployment under strict on-device resource constraints. In telecommunication fraud scenarios, a typical scam call lasts approximately 90~s, while critical fraudulent instructions are often delivered within the first 60~s, motivating a strict latency-aware system design.

To support real-time intervention, we introduce two distinct operational latency metrics:
\begin{enumerate}
    \item \textit{Per-window incremental inference latency} refers to the time required for each risk-scoring iteration triggered by a sliding temporal window. We constrain this latency to $\le 500$~ms to enable continuous updates at a 1-s interval; our system achieves a median ($P_{50}$) latency of 268~ms (see Section~\ref{sec:sysarch}).
    \item \textit{First-alert latency} measures the time from call initiation to the first risk prediction. Given a 3-s acoustic window ($W$), the system produces an initial alert within 3.3~s, enabling early-stage intervention during the coercion phase.
\end{enumerate}

Beyond the privacy mandates of PIPL, the application side is subject to tight hardware constraints---a 4 GB local memory footprint and 500 MB of available storage---which jointly motivate efficient on-device inference, wherein 4-bit quantised Large Language Models (LLMs) offer a pragmatic trade-off among computational efficiency, privacy preservation, and real-time responsiveness.



In our empirical evaluations, standard post-training quantisation (PTQ) at the 0.5B scale leaves an $F_1$-score gap of 7.1--8.5 percentage points relative to our quantisation-aware distillation (QAD) framework on \textit{TAF-28k} (Section~\ref{sec:main-results}, Table~\ref{tab3-en}). Conventional quantisation-aware training (QAT), by contrast, re-establishes the full supervised fine-tuning (SFT) and reinforcement-learning (RL) pipelines, which inflates engineering complexity and risks the alignment properties acquired during initial model training. The NVIDIA Nemotron report~\citep{3} indicates that KL-divergence-based distillation can recover approximately 97.0\%--99.4\% of full-precision performance; we treat these findings as a reference baseline and independently validate them on the target anti-fraud task (Section~\ref{sec:main-results}).

Directly transferring the native QAD paradigm to telecommunication fraud detection raises three practical challenges:

\paragraph{(C1) Privacy-constrained data and acoustic processing.}
China's PIPL restricts both the cross-border transmission of raw audio/biometric data and the sharing of fraud-call corpora, which jointly limits the availability of high-quality labelled multimodal data and motivates on-device acoustic processing. This dual constraint demands a reconstruction-resistant acoustic representation that resists speech-content reconstruction while remaining sufficiently informative for downstream multimodal fusion; the resulting on-device embedding $\bm{F}_v$ is precisely what lets raw biometric audio stay within the device boundary motivated by PIPL's data-minimisation principle.

\paragraph{(C2) Robustness under ultra-low-bit quantisation.}
Low-bit quantisation formats such as \texttt{NVFP4} and \texttt{GGUF} degrade accuracy substantially, yet existing QAD investigations focus on single-modality language tasks and do not address feature-distribution alignment under these formats in fraud-detection settings; deployed baselines such as SAFE-QAQ likewise lack deployment compatibility under ultra-low-bit formats.

\paragraph{(C3) Real-time inference on resource-constrained edge devices.}
Risk scores must be refreshed at continuous 1-s intervals, requiring the per-window inference latency to remain strictly below 500~ms (see Section~\ref{sec:streaming}). Although low-bit quantisation inherently reduces computational overhead, further domain-specific optimisation is required to satisfy such stringent latency constraints on commodity mobile hardware, motivating the integration of speculative decoding mechanisms.

To address these three challenges, we propose \textit{QAD-MultiGuard}, an edge--cloud collaborative framework for multimodal fraud detection and risk assessment. The overarching objective is to deliver, on resource-constrained mobile hardware, a multimodal fraud detector that simultaneously satisfies three coupled requirements: (i)~\textit{privacy}---raw audio and text never leave the device; (ii)~\textit{fidelity}---the 4-bit quantised student recovers the BF16 teacher's accuracy; and (iii)~\textit{responsiveness}---per-window inference remains below 500~ms. Corresponding to these requirements, the framework extends the native NVFP4 QAD paradigm along four dimensions, each addressing a specific challenge:

\paragraph{(1) Pure-KL homologous distillation and dual-track quantisation (addressing C2).}
A pure KL-divergence distillation strategy is adopted to transfer structural knowledge from a full-precision BF16 teacher to heterogeneous low-bit student architectures, thereby recovering the accuracy degraded by weight quantisation. A dual-track quantisation design seamlessly pairs high-throughput cloud-side NVFP4 deployment with resource-constrained on-device GGUF-Q4\_K\_M inference.

\paragraph{(2) Output-Variance Freeze (OV-Freeze) regularisation (addressing C2).}
An OV-Freeze regularisation strategy is introduced to further stabilise quantisation-sensitive projection layers by dynamically aligning student output statistics with those of the teacher model under ultra-low-bit conditions, preventing representation collapse and complementing the distillation objective of (1).

\paragraph{(3) Reconstruction-resistant acoustic representation and multimodal fusion (addressing C1).}
A compact 128-dimensional reconstruction-resistant acoustic representation is designed by combining handcrafted and deep neural speech features, enabling robust audio--text fusion while substantially elevating reconstruction errors under adversarial speech-content attacks, thereby allowing raw audio to remain on-device. This representation provides linguistic/semantic content privacy---resisting speech-content reconstruction (WER $\ge 0.95$, exceeding the $\ge 0.90$ threshold of the threat model in Section~\ref{Formalised_Threat_Model})---but is not designed for cross-session unlinkability; a semi-honest cloud may still link repeated sessions of the same speaker through embedding similarity, a residual risk analysed in Section~\ref{sec:discussion}. The decision-fusion framework is four-modal (text, acoustic, URL, and metadata); the empirical evaluation is exercised on the text and acoustic modality pair provided by TAF-28k, with the URL and metadata fusion weights carried as deployment parameters.

\paragraph{(4) Domain-adapted speculative decoding and edge--cloud deployment (addressing C3).}
A domain-specialised speculative decoding framework accelerates the asynchronous cloud-side chain-of-thought review, while the on-device critical path meets the sub-500~ms latency budget through a quantised student and lightweight linear fusion. Empirical results show that the on-device \texttt{Q4\_K\_M} student attains an $F_1$-score of 0.917 with a median ($P_{50}$) latency of 268~ms on the \textit{TAF-28k} corpus, and the cloud-side \texttt{NVFP4} track reaches 0.923 under asynchronous review, while maintaining structural robustness across multiple cross-domain datasets.

These four components are carried by a three-tier edge--cloud architecture (Section~\ref{sec:sysarch}) whose on-device decision-level fusion layer (Tier-3) integrates the modality risk scores into a single real-time decision, constituting the deployment skeleton for the four components.

The principal contributions of this work are threefold. \textit{Methodologically}, we introduce a pure-KL homologous distillation objective that recovers the accuracy lost to 4-bit quantisation without the representation-space mismatch of cross-scale distillation, complemented by an Output-Variance Freeze (OV-Freeze) regulariser that stabilises quantisation-sensitive projection layers; the measured effect of OV-Freeze is best interpreted as evidence of faithful distillation rather than a large accuracy improvement (Section~\ref{sec:main-results}). \textit{On privacy}, we construct a 128-dimensional acoustic embedding that keeps raw biometric audio on-device while preserving the coarse prosodic cues necessary for fraud detection, advancing representation-level protection beyond ASR-transcript pipelines. \textit{On systems}, we co-design an edge--cloud deployment in which a quantised on-device student meets the sub-500~ms latency budget while an asynchronous cloud review---accelerated by domain-adapted speculative decoding---refines the decision off the critical path. The novelty lies not in any single component---each draws on established techniques---but in their integrated co-design and empirical validation for privacy-constrained on-device fraud detection, a multimodal setting that existing single-modality and server-only baselines do not address.

\vspace{0.5em}
We note upfront that the main multimodal evaluation relies on a single primary corpus (\textit{TAF-28k}), collected under a controlled re-enactment protocol; complementary adversarial (\textit{AdvFraud-3k}) and cross-domain (\textit{ChiFraud}) assessments are reported, but generalisation to unconstrained real-world conditions is not yet formally validated (see Section~\ref{sec:discussion}).

The remainder of this paper is organised as follows. Section~\ref{sec:related} reviews the related literature. Section~\ref{sec:method} details the system architecture and proposed methodology. Section~\ref{sec:experiments} reports the experimental evaluations and comparative results. Section~\ref{sec:discussion} discusses inherent limitations, and Section~\ref{sec:conclusion} concludes the paper.

\section{Related Work}\label{sec:related}
\subsection{Quantisation-Aware Distillation (QAD)}

The knowledge-distillation paradigm introduced by \citet{4} transfers structural knowledge from a high-capacity teacher network to a compact student network via soft-label supervision. \citet{5} first integrated knowledge distillation (KD) into weight quantisation, establishing quantised distillation as a distinct trajectory. \citet{6} demonstrated that low-precision networks can effectively fit the output representations of full-precision teachers to markedly improve downstream accuracy. Furthermore, \citet{7} proposed LLM-QAT, validating data-free quantisation-aware training utilising model-generated calibration data, whilst \citet{8} introduced BitDistiller, combining asymmetric quantisation with self-distillation objectives to mitigate representation degradation in sub-4-bit Large Language Models (LLMs).

Recent advances in LLM quantisation can be broadly divided into distillation-based approaches and post-training quantisation (PTQ) methods. The NVIDIA Nemotron technical report~\citep{3} suggests that, relative to conventional QAT, QAD may better preserve the capabilities of alignment-tuned models by avoiding hard-label supervision during quantised training. The report further indicates that homologous-teacher alignment---wherein the BF16 teacher and student share the identical architecture and pre-training lineage---and KL-divergence objectives are particularly effective for low-bit distillation.

On the PTQ side, several frameworks have achieved substantial progress without requiring an intensive distillation phase. \citet{9} proposed activation-aware weight quantisation (AWQ), which protects critical channel weights via activation-aware scaling. \citet{10} developed GPTQ, which leverages second-order Taylor information for layer-wise weight reconstruction. Additionally, \citet{11} introduced SpinQuant, which suppresses activation outliers via learnable orthogonal rotations, \citet{12} presented QuaRot, which achieves outlier-free 4-bit inference via Hadamard transforms, \citet{xiao2023smoothquant} proposed SmoothQuant, which migrates the quantisation difficulty from activations to weights via per-channel smoothing, \citet{shao2024omniquant} introduced OmniQuant, which calibrates weight and activation quantisation omnidirectionally, and \citet{dettmers2023qlora} combined 4-bit quantisation with low-rank adapters (QLoRA) to finetune quantised models without restoring full-precision weights. Although these PTQ methodologies have demonstrated competitive performance on general benchmarks, our empirical evaluations reveal a severe performance gap of 7.1--8.5 percentage points on the domain-specific \textit{TAF-28k} benchmark under 4-bit deployment configurations (see Section~\ref{sec:main-results}), because weight-only PTQ preserves the magnitude statistics of pre-trained weights but cannot recover the domain-specific cross-modal features that fraud detection relies on, leaving them vulnerable to the distributional shift induced by low-bit reconstruction.

However, existing QAD investigations have largely focused on general-purpose language modelling and have not yet been systematically evaluated in privacy-sensitive, multimodal fraud-detection settings. Accordingly, this work extends the homologous-teacher QAD paradigm to this underexplored territory of privacy-sensitive multimodal telecommunication-fraud analysis, and introduces an Output-Variance Freeze (OV-Freeze) regulariser to enhance architectural robustness under ultra-low-bit quantisation formats.

\subsection{Telecommunication Fraud Detection}

Early anti-fraud detection systems heavily relied on keyword-matching rule engines or traditional machine learning (ML) classifiers, such as support vector machines (SVMs) and Random Forests (RF)~\citep{13,14,15}. Recent telecommunication fraud detection methodologies can be broadly categorised into three dominant streams: text-based approaches, graph-based relational reasoning frameworks, and multimodal detection systems.

Text-based methods have extensively leveraged pre-trained language models to elevate detection accuracy. For instance, BERT-based fraudulent-message classification via sentiment analysis~\citep{16} has been explored for scam-message screening, whilst RoBERTa-MHARC~\citep{liRobertaMHARC2024} integrated RoBERTa with a multi-head attention mechanism and a dual-loss objective on a five-category telecommunication fraud corpus. More recently, an LLM-based multi-role, multi-layer prompting strategy~\citep{mrml2026} introduced hierarchical reasoning to resolve long-tail and ambiguous fraud variants. Additionally, attention-based convolutional networks applied to communication-feature images~\citep{liImageCNN2025} have been developed to address the severe class imbalances inherent to operational call data.

Graph-based relational reasoning over caller-device-metadata networks has also been explored~\citep{papasavva2025,MotieESWA2024,RenESWA2024}, but remains largely confined to network-level behavioural features rather than the privacy-constrained on-device inference targeted in this work.

Despite these methodological advances, existing approaches remain largely confined to isolated textual or behavioural features, and none of them address quantisation-aware distillation for low-bit edge inference; they also lack robust architectural support for on-device deployment and fail to address privacy-constrained acoustic-modality fusion. These are precisely the paradigm gaps addressed by \textit{QAD-MultiGuard}.

To provide empirical benchmarks, \citet{1} released \textit{TeleAntiFraud-28k} (\textit{TAF-28k}), the first open-source audio--text dual-modality dataset featuring complex reasoning annotations specialized for the telecommunication anti-fraud domain. The corpus contains 28,511 anonymised voice--text pairs, totalling over $307\text{~h}$ of acoustic data. Subsequently, \citet{2} proposed \textit{SAFE-QAQ}, an RL-based end-to-end audio--text multimodal framework tailored exclusively for high-performance server-side deployment; it does not accommodate resource-constrained mobile environments.

The proposed \textit{QAD-MultiGuard} framework builds upon these foundational advances while executing a distinct operational objective. Whereas \textit{TAF-28k} establishes a rigorous benchmark for cross-modal reasoning and \textit{SAFE-QAQ} substantiates the utility of end-to-end multimodal architectures, the present work uniquely concentrates on privacy-constrained on-device inference via ultra-low-bit quantisation within an edge--cloud collaborative deployment paradigm. Table~\ref{tab1} systematically summarises the core structural distinctions between existing baseline systems and our proposed framework.

\begin{table*}[t]
\centering
\caption{Systematic comparison with representative telecommunication fraud detection studies.}
\label{tab1}
\renewcommand{\arraystretch}{1.15}
\setlength{\tabcolsep}{8pt}
\begin{tabular}{@{} m{3cm} m{3.5cm} m{3.2cm} m{5cm} @{}}
\toprule
\textbf{Dimension} 
& \textbf{TAF-28k\newline\citep{1}} &\textbf{SAFE-QAQ\newline\citep{2}} 
& \textbf{QAD-MultiGuard (This work)}\\
\midrule
Core role 
& Dataset \& benchmark & Multimodal framework & Edge--cloud framework \\
\addlinespace[2pt]
Audio representation & ASR-derived & Learned & Reconstruction-resistant $128\text{-d}$ embedding \\
\addlinespace[2pt] Quantisation support & N/A & N/A & NVFP4 $+$ GGUF-Q4\_K\_M \\
\addlinespace[2pt] Deployment paradigm  & N/A & Centralised inference  & Edge--cloud collaborative inference \\
\addlinespace[2pt]
Privacy mechanism 
& Data anonymisation 
& Not reported 
& Representation-level protection \\
\addlinespace[2pt]
Training paradigm 
& SFT baseline 
& RL-based optimisation 
& Pure-KL homologous distillation \\
\addlinespace[2pt]
Empirical validation 
& Benchmark verification 
& Production deployment& Benchmark + on-device evaluations \\
\bottomrule
\end{tabular}
\end{table*}

\subsection{On-device Large Language Model Inference}

On-device inference of Large Language Models (LLMs) is fundamentally constrained by model storage footprints, local memory bandwidth, and sequential decoding latency. The \texttt{GGUF} format~\citep{17} compresses 0.5B-parameter architectures to approximately $240\text{~MB}$ under Q4\_K\_M mixed-precision block-wise quantisation configurations, thereby enabling local deployment on mobile ARM-based systems-on-chip (SoCs). Recent empirical investigations have demonstrated the engineering feasibility of deploying such compressed architectures on resource-constrained edge hardware for real-time diagnostic applications~\citep{PereiraESWA2025}.

To alleviate the hardware latency bottlenecks inherent to autoregressive decoding, \citet{20} introduced speculative decoding, wherein a lightweight draft model sequentially generates multiple candidate tokens that are subsequently verified in parallel by the primary target model. \citet{21} formally derived the analytical relationship between the average token-acceptance rate $\alpha$ and the resulting sequence-generation computational acceleration:
\begin{equation}
\label{eq:speedup}
\text{Speedup} = \frac{1-\alpha^{\gamma+1}}{1-\alpha},
\end{equation}
where $\gamma$ denotes the fixed number of candidate tokens generated per draft iteration step. Rather than employing a generic cross-domain draft model, this work fine-tunes the compact draft architecture specifically on target domain-specific anti-fraud corpora to precisely capture the recurrent linguistic structures of fraud-engineering scripts. This domain-adapted alignment successfully increases the average token-acceptance rate $\alpha$ from 0.78 to 0.86, substantially lowering end-to-end response times.

\subsection{Reconstruction-Resistant Acoustic Representations}

China's Personal Information Protection Law (PIPL) imposes strict restrictions on the external transmission of raw biometric information, mandating that such sensitive data remain locally on user terminals unless explicit authorisation has been granted and a rigorous prior security assessment has been conducted. \citet{snyder2018xvectors} introduced x-vectors for speaker recognition, which were subsequently adapted as an early approach for acoustic identity obfuscation. However, \citet{23} demonstrated, utilising the Generative Latent Optimisation (GLO) attack framework, that such anonymised embeddings can still be exploited for inverse reconstruction. Separately, anti-spoofing and deepfake-speech detection have been standardised through the ASVspoof challenge series~\citep{yamagishi2021asvspoof}, which evaluates countermeasures against synthesised, converted, and replayed speech; this detection-oriented line of work complements, rather than replaces, the representation-level reconstruction resistance pursued here. More recent speaker-anonymisation research has progressed from heuristic obfuscation toward formal guarantees---including $k$-anonymity in embedding space and differentially-private perturbation of the representation---and is standardly evaluated through the Voice Privacy Challenge framework~\citep{22}, which distinguishes \textit{content} from \textit{identity} dimensions: the content dimension is evaluated as a utility objective (anonymised speech should preserve linguistic content, i.e.\ low ASR word-error rate), whereas the identity dimension targets cross-session unlinkability of the speaker. The scope of the present work is deliberately \textit{content-private} in a stronger sense than VPC's content-utility: our representation resists speech-content reconstruction (high WER) rather than preserving content, and is not designed for identity-level unlinkability, a boundary analysed in Section~\ref{sec:discussion}.

To address this issue, we design a compact $128\text{-dimensional}$ acoustic representation that reduces reconstructable information in the feature space. Specifically, time-averaged Mel-frequency cepstral coefficient (MFCC) features and globally average-pooled Whisper-tiny encoder features are concatenated to form a $128\text{-dimensional}$ embedding $\bm{F}_v$, specified as a design target; the released experiments evaluate these two components through their respective proxy embeddings (Section~\ref{sec:acoustic}). This aggregation effectively reduces fine-grained frame-level temporal structure, resulting in a Word Error Rate (WER) of $\ge 0.95$ under both GLO and model-inversion adversarial attack settings. Further technical details and empirical validations are provided in Section~\ref{sec:acoustic} and Appendix~\ref{app:nonrecon}, respectively.

The overall system architecture is presented in Section~\ref{sec:method}, covering four components---quantisation-aware distillation (QAD), Output-Variance Freeze (OV-Freeze), reconstruction-resistant acoustic representations, and domain-adapted speculative decoding---together with the decision-level fusion layer that integrates them.

\section{System Architecture and Method}\label{sec:method}
\subsection{System Architecture and Threat Model}
\subsubsection{System Architecture}\label{sec:sysarch}

QAD-MultiGuard adopts a three-tier edge--cloud defence architecture (Figure~\ref{fig1}):

\begin{figure}
    \centering
    \includegraphics[width=\linewidth]{figure/fig1_architecture.png}
    \caption{\textit{QAD-MultiGuard} three-tier edge--cloud collaborative architecture. Raw audio, ASR transcripts, and SMS plaintext remain strictly on-device. Only the $128\text{-dimensional}$ reconstruction-resistant acoustic representation $\bm{F}_v$ and scalar risk scores are transmitted to the cloud infrastructure. In branch (c), only the Whisper-tiny encoder is deployed on-device, whilst the decoder module is completely omitted and no ASR transcript is generated. The $(\epsilon,\delta)$-local differential privacy ($(\epsilon,\delta)$-LDP, $\epsilon=1.5$, $\delta=10^{-5}$) mechanism is included as an optional configuration for enhanced-privacy scenarios and is omitted from the main experimental configuration (see Section~\ref{sec:discussion}).}
    \label{fig1}
\end{figure}

\paragraph{Tier-1: On-device real-time detection.}
Deployed on Snapdragon 8 Gen 3 and comparable mobile platforms, this tier executes a lightweight student architecture ($240\text{~MB}$, \texttt{Q4\_K\_M} quantised), which incorporates an acoustic embedding module together with SMS ($12\text{-d}$) and URL ($6\text{-d}$) feature extractors supporting the four-modal fusion framework. The acoustic feature vector $\bm{F}_v \in \mathbb{R}^{128}$ is computed locally, and the student model produces text and acoustic risk scores that are fused on-device (Tier-3) into an immediate risk decision. This tier constitutes the blocking real-time path: raw audio streams and SMS plaintexts never leave the device, thereby aligning with data-minimisation practices without incurring any cloud round-trip latency.

\paragraph{Data-flow boundary.}
For voice interactions, the Whisper-tiny encoder is executed locally without its corresponding decoder component, ensuring no ASR transcript is generated on-device. For SMS traffic, the local student model processes the $12\text{-d}$ structural features to output a scalar text risk score $r_{\text{text}} \in [0,1]$. Along the asynchronous review path, the cloud-side tier receives exclusively: (i) the $128\text{-dimensional}$ embedding $\bm{F}_v$, (ii) the text and acoustic risk scores, and (iii) intermediate hidden states intended for cross-modal consistency verification. Raw audio, ASR text transcripts, and raw SMS plaintexts remain strictly confined within the local device boundary.

\paragraph{Deployment model.}
The on-device privacy boundary presupposes a deployment model in which the detection software executes on the data subject's own mobile device---the protected call participant---rather than within the operator's network. This work adopts that model as an explicit design assumption: the ``raw audio stays on-device'' property holds only when the framework runs on the victim's handset. An alternative operator-side deployment, in which the operator already processes audio in-network under telecommunications regulation, implies a different privacy boundary that the present on-device architecture does not cover (Section~\ref{sec:discussion}). No legal-compliance claim is made with respect to either deployment model.

\paragraph{Tier-2: Cloud asynchronous review.}
Targeted for deployment on an NVIDIA Blackwell GPU cluster, this tier operates an \texttt{NVFP4}-quantised 0.5B-parameter LLM dedicated to chain-of-thought reasoning and comprehensive cross-modal consistency verification. The cloud-side architecture processes intermediate hidden states transmitted from the edge layer, which share a homologous structural lineage with the on-device student model. Crucially, Tier-2 operates asynchronously and does not block the on-device decision path: Tier-1 and Tier-3 issue the immediate risk decision locally, while the cloud review returns a refined confirmation in parallel. Review requests are scheduled by risk level so that medium-risk scenarios ($0.35 \le r < 0.70$), which benefit most from deep reasoning, receive priority, whereas unambiguous low- and high-risk cases are skipped or deprioritised. Should the asynchronous review overturn an on-device decision---for example, reclassifying a locally low-risk call as high-risk after deeper reasoning---a secondary alert is issued through the same Tier-3 fusion path, so that a divergent cloud verdict is never silently discarded.

\paragraph{Tier-3: Multimodal risk fusion.}
This tier is executed locally on-device and performs decision-level fusion across the text, acoustic, URL, and metadata modality risk scores utilising a weighted sigmoid aggregation scheme. On TAF-28k, which supplies only text and acoustic scores, the text and acoustic fusion weights are optimised offline via the L-BFGS algorithm, while the URL and metadata weights are carried forward as deployment parameters (Eq.~\eqref{eq:w-deploy}), enabling deterministic, real-time three-level risk decisions.

\subsubsection{Formalised Threat Model}\label{Formalised_Threat_Model}

We consider three representative threat vectors: $\mathcal{G}_1$ (privacy reconstruction attacks), $\mathcal{G}_2$ (adversarial evasion attacks, including acoustic jamming interventions such as \texttt{ASRJam}~\citep{19}), and $\mathcal{G}_3$ (identity fraud spoofing attacks; addressed at the design level via cross-modal consistency verification, and not quantitatively benchmarked in this study).

The corresponding security and robustness constraints are strictly specified as follows:
\begin{enumerate}
    \item[\text{(i)}] Under privacy reconstruction threat $\mathcal{G}_1$, acoustic reconstruction quality is lower-bounded by an adversary Word Error Rate (WER) threshold of $\ge 0.90$; the deployed embedding empirically achieves $\text{WER} \ge 0.95$ (Section~\ref{sec:glo}), providing strong resistance to reconstruction under the evaluated attack models.
    \item[\text{(ii)}] Under evasion threat $\mathcal{G}_2$, the relative degradation in $F_1$-score, evaluated against a matched full-precision BF16 baseline, is constrained to be $\le 6\%$.
    \item[\text{(iii)}] Under identity fraud threat $\mathcal{G}_3$, the system enforces strict cross-modal consistency verification across acoustic, textual, and metadata-derived risk signals to mitigate multimodal exploitation.
\end{enumerate}

\subsection{Quantisation-Aware Distillation (QAD)}
\subsubsection{Pure KL-Divergence Loss}

Given an input sequence $x$ and a vocabulary $\mathcal{V}$, let $p_{\text{teacher}}(y|x)$ and $p_{\text{student}}(y|x)$ denote the token-level output probability distributions of the full-precision BF16 teacher model and the quantised student model, respectively. The primary QAD optimisation objective is defined as:
\begin{equation}
\label{eq:kl-loss}
L_{\mathrm{QAD}} = D_{\mathrm{KL}}\!\bigl( p_{\text{teacher}}(y|x)\,\|\,p_{\text{student}}(y|x) \bigr).
\end{equation}

The softmax temperature is strictly fixed to $T = 1$ during distillation to guarantee distribution consistency between the training and inference phases. In contrast to conventional quantisation-aware training (QAT) frameworks~\citep{24,25,26}, which retain a hard cross-entropy term inherited from supervised fine-tuning (SFT), the hybrid QAT objective inherently introduces an artificial distributional mismatch between teacher and student outputs.

Empirically, our proposed pure-KL formulation achieves a consistently lower divergence from the native BF16 teacher distribution ($D_{\mathrm{KL}} = 0.005$) compared with the hybrid QAT baseline ($D_{\mathrm{KL}} = 0.311$). Detailed ablation studies validating this loss formulation are provided in Section~\ref{sec:loss-ablation} (see Table~\ref{tab5-en}).

\subsubsection{Homologous teacher self-distillation}

The native Qwen2.5-0.5B-Instruct model (BF16)~\citep{27,28,29} is adopted as the homologous teacher. QAD is designed to correct output-distribution shifts induced by low-bit quantisation rather than to transfer task-specific knowledge from heterogeneous teacher architectures.

\subsubsection{Edge--Cloud Co-Quantisation}

We propose an edge--cloud dual-track quantisation framework wherein a server-side model and an on-device model are deployed utilising \texttt{NVFP4} (Blackwell GPUs) and \texttt{Q4\_K\_M} (ARM/Snapdragon platforms), respectively. Both architectures share an identical full-precision BF16 backbone and are independently optimised through quantisation-aware distillation using the proposed pure-KL objective. This decoupled design enables strict distribution alignment at the token output level without requiring explicit, computationally expensive feature-space regularisation across distinct deployment tiers.

On the server side, \texttt{NVFP4} training is implemented via quantise--dequantise (QDQ) fake-quantisation operators on H100-class hardware, while the on-device student model undergoes \texttt{GGUF}-style block-wise fixed-point quantisation. Although these two quantisation schemes diverge in numerical representation and hardware execution paths, both are bounded by the exact same BF16 teacher distribution, which provides a unified optimisation target across heterogeneous deployment environments.

Table~\ref{tab2-en} systematically summarises the precise structural and operational configurations of this dual-track quantisation paradigm.

\begin{table*}[t]
\centering
\renewcommand{\arraystretch}{1.15}
\caption{Systematic configuration of the edge--cloud dual-track quantisation scheme.}
\label{tab2-en}
\setlength{\tabcolsep}{8pt}
\begin{tabularx}{\textwidth}{>{\raggedright\arraybackslash}p{0.25\linewidth}XX}
\toprule
\textbf{Component Dimension} & \textbf{Server Track (\texttt{NVFP4})} & \textbf{On-Device Track (\texttt{Q4\_K\_M})} \\
\midrule
Backbone Architecture & \texttt{Qwen2.5-0.5B-Instruct} & \texttt{Qwen2.5-0.5B-Instruct} \\
Active Parameter Scale & 494M & 494M \\
Distillation Teacher & BF16 Homologous Self-Distillation & BF16 Homologous Self-Distillation \\
Quantisation Format & \texttt{NVFP4} (block size = 16, FP8 E4M3 scaling) & \texttt{Q4\_K\_M} (\texttt{GGUF} block quantisation) \\
Compressed Model Footprint & $248\text{~MB}$ & $240\text{~MB}$ \\
Target Deployment Platform & NVIDIA Blackwell GPU & Mobile ARM / Snapdragon SoC \\
Hardware-Level Acceleration & $2.1\times$ (kernel-level throughput vs FP8) & --- \\
Accuracy Recovery vs BF16 Baseline & 99.1\% & 98.5\% \\
\bottomrule
\end{tabularx}

\vspace{1mm}
\footnotesize
Note: the ``NVIDIA Blackwell GPU'' entry denotes the target deployment platform; all \texttt{NVFP4} accuracy results reported in this work are generated via the QDQ numerical-behaviour-emulation (NBE) protocol on H100 hardware (Section~\ref{sec:nbe}), not native Blackwell execution. The compressed-model footprints denote weight-only sizes; the on-disk \texttt{GGUF} container for the \texttt{Q4\_K\_M} student is larger ($\approx 491\text{~MB}$).
\end{table*}

\paragraph{Heterogeneous quantisation evaluation.}
The empirical effectiveness of this asymmetric dual-track design is quantified in Section~\ref{sec:quant-ablation} (see the quantisation-scheme ablation of Figure~\ref{fig4}a), where the heterogeneous configuration is benchmarked against a homogeneous INT4-style baseline under an otherwise identical pipeline.

\paragraph{Cross-tier interface design.}
The cross-tier communication boundary handles exclusively FP16-formatted intermediate hidden representations. Low-bit edge tensors are locally dequantised prior to packet transmission across the network boundary to enforce privacy containment (Eq.~\ref{eq:dequant}). Specifically, the edge terminal transmits: (i) the reconstruction-resistant $128\text{-dimensional}$ acoustic embedding $\bm{F}_v$, (ii) the text and acoustic risk scores, and (iii) selected intermediate hidden states reserved for cross-modal consistency cross-checking. Raw biometric voice data, text payload strings, and ASR transcripts remain entirely confined within the local hardware boundary.

Let $\bm{h}_{\mathrm{int}}$ denote a quantised hidden representation block with an associated scale factor $s_{\mathrm{blk}}$ and zero-point parameter $z$. The runtime dequantisation operation is mathematically formulated as:
\begin{equation}
\label{eq:dequant}
\tilde{\bm{h}} = (\bm{h}_{\mathrm{int}} - z) \cdot s_{\mathrm{blk}} \to \text{FP16}.
\end{equation}

Because both endpoints trace their architectural lineage back to the same homologous BF16 teacher, the cross-tier activation distributions remain closely aligned, requiring minimal statistical adjustment during cross-tier handshakes.

\paragraph{System-level interpretation.}
The reported $2.1\times$ hardware speedup reflects the isolated compute-kernel acceleration of \texttt{NVFP4} dense matrix operations relative to standard native FP8 operations on Blackwell architectures. This gain therefore represents a raw hardware execution throughput margin rather than an equivalent reduction in end-to-end user latency, as the overall operational response cycle remains governed by local feature extraction, encrypted packet transmission overheads, and multimodal decision-fusion layers.

\subsubsection{Numerical Behaviour Emulation (NBE) Protocol}\label{sec:nbe}

A Numerical Behaviour Emulation (NBE) protocol~\citep{NVIDIA2024TRTLLMReleaseNotes,NVIDIA2024TensorRTLLMDocs} is implemented on H100 hardware infrastructures to guarantee a rigorous, consistent comparison between compressed low-bit execution paths and full-precision baselines. The underlying protocol systematically embeds quantise--dequantise (QDQ) fake-quantisation operators directly within the forward computational graph:
\begin{equation}
\label{eq:nbe}
\widehat{\bm{W}} = \text{clamp}\!\bigl(\text{round}(\bm{W}/s),\,q_{\min},\,q_{\max}\bigr) \cdot s,
\end{equation}
where the composite dual-layer scale factor is defined as $s = s_{\text{block}} \cdot s_{\text{tensor}}$, with the block-wise scale represented in the FP8 E4M3 format (Table~\ref{tab2-en}). Eq.~\eqref{eq:nbe} thus states the general integer round-and-clamp form of the QDQ operators: under the emulation protocol the weight grid is an integer grid, distinct from the native FP4 (E2M1) grid of hardware \texttt{NVFP4}, a distinction made explicit for transparency. Empirical cross-comparison evaluations benchmarking this emulation setup against the official NVIDIA \texttt{NVFP4} reference checkpoint across \textit{MMLU}~\citep{30}, \textit{GSM8K}~\citep{31}, and \textit{HumanEval}~\citep{32} yield an absolute evaluation accuracy deviation of $< 0.3\%$, mathematically verifying the numerical correctness of the integrated QDQ pipeline. Consequently, all \texttt{NVFP4} results reported in this work---including the main comparisons of Table~\ref{tab3-en} and the cross-dataset evaluations of Table~\ref{tab4-cross-en}---are generated strictly utilising this standardised emulation protocol.

\subsubsection{Output-Variance Freeze (OV-Freeze) Regularisation}\label{sec:ovfreeze-method}

OV-Freeze enforces strict variance consistency between the quantised student network and the full-precision BF16 teacher by aligning layer-wise activation statistics under a calibrated reference distribution.

The regularisation objective is defined as:
\begin{equation}
\label{eq:ovf-loss}
L_{\mathrm{OVF}} = \lambda \sum_{\ell \in \mathcal{P}} \left\| \operatorname{Var}^{(t)}_{\mathrm{EMA}}(\bm{y}_\ell) - \sigma^2_{\mathrm{BF16},\ell} \right\|_2^2,
\end{equation}
where $\mathcal{P} = \{q, k, v, o\}_{\text{proj}}$ denotes the set of target attention projection layers, $\sigma^2_{\mathrm{BF16},\ell}$ represents the calibrated variance statistic derived from static BF16 baseline inference, and $\lambda = 0.01$ is a scaling hyperparameter. A detailed mathematical stability analysis of the projection-layer rescaling, including the boundedness of the correction factor and its effect on optimisation stability, is provided in Appendix~\ref{app:bounded}.

The variance estimator is dynamically updated via an exponential moving average (EMA):
\begin{equation}
\label{eq:ema}
\operatorname{Var}^{(t)}_{\mathrm{EMA}}(\bm{y}_\ell)
= \rho \cdot \operatorname{Var}^{(t-1)}_{\mathrm{EMA}}(\bm{y}_\ell)
+ (1-\rho) \cdot \operatorname{Var}^{(t)}_{\mathrm{batch}}(\bm{y}_\ell),
\quad \rho = 0.95,
\end{equation}
which effectively stabilises second-order activation statistics under stochastic optimisation trajectories.

The unified joint optimisation objective is formulated as:
\begin{equation}
\label{eq:joint}
L_{\mathrm{joint}} = L_{\mathrm{QAD}} + L_{\mathrm{OVF}},
\end{equation}
where $L_{\mathrm{QAD}}$ corresponds to the primary distribution distillation loss defined in Eq.~\eqref{eq:kl-loss}.

During the forward pass, a variance-aware rescaling transformation is dynamically applied to the layer outputs:
\begin{equation}
\label{eq:ovf-rescale-en}
\bm{y}'_\ell = \bm{y}_\ell \cdot c_\ell, \qquad
c_\ell = \operatorname{sg}\!\left[\sqrt{\frac{\sigma^2_{\mathrm{BF16},\ell}}{\operatorname{Var}^{(t)}_{\mathrm{EMA}}(\bm{y}_\ell) + \epsilon}}\,\right].
\end{equation}

The introduction of the stop-gradient operator ($\operatorname{sg}$) yields the following simplified backward error propagation flow:
\begin{equation}
\label{eq:ovf-detach-en}
\frac{\partial L_{\mathrm{QAD}}}{\partial\bm{y}_\ell}
= c_\ell \cdot \frac{\partial L_{\mathrm{QAD}}}{\partial\bm{y}'_\ell}.
\end{equation}

The forward rescaling of Eq.~\eqref{eq:ovf-rescale-en} and the variance-matching loss of Eq.~\eqref{eq:ovf-loss} in the joint objective (Eq.~\eqref{eq:joint}) are complementary rather than redundant: the stop-gradient factor $c_\ell$ performs instantaneous numerical alignment of activations during the forward pass, whereas $L_{\mathrm{OVF}}$ steers the learnable parameters---through gradient back-propagation---toward regions where the variance gap is intrinsically small.

Operationally, the OV-Freeze module is activated exclusively during the final 30\% of the training schedule. OV-Freeze is a training-time regulariser only: the correction factor $c_\ell$ is applied solely during this terminal phase and is removed at evaluation and inference, which therefore measure the plain quantised student. During this terminal phase, the student fully inherits the teacher's attention mechanism by aligning the output variance of the $\{q,k,v,o\}$ projection layers, and freezing this variance over the last 30\% of training prevents representational collapse during the distillation-to-deployment transfer. Empirically, this targeted regularisation consistently improves convergence stability and mitigates projection-layer variance drift. Under an ultra-low-bit \texttt{Q4\_K\_M} quantisation profile, the layer-wise variance deviation decreases from $+18.2\%$ down to $+1.3\%$ relative to the native BF16 reference baseline statistics.

\begin{figure}
\centering
\includegraphics[width=\linewidth]{figure/fig2_acoustic_embedding.pdf}
\caption{Schematic construction of the $128\text{-dimensional}$ reconstruction-resistant acoustic embedding $\bm{F}_v$. (a)~Log-Mel spectrogram representation; (b)~time-averaged MFCC representation; (c)~concatenated joint embedding $\bm{F}_v \in \mathbb{R}^{128}$ constructed from the local MFCC features and globally pooled \texttt{Whisper-tiny} encoder representations. Dual temporal aggregation compresses frame-level fine structures into compact, reconstruction-resistant statistics for downstream secure inference.}
\label{fig2}
\end{figure}

\subsection{Reconstruction-Resistant Acoustic Embedding}\label{sec:acoustic}

The reconstruction-resistant acoustic embedding is specified as a $128\text{-dimensional}$ design target:
\begin{equation}
\label{eq:f-v}
\bm{F}_v = \bigl[\bm{f}_{\mathrm{mfcc}};\,\psi(\bm{W}_{\mathrm{proj}}\bar{\bm{h}}_w)\bigr]\in\mathbb{R}^{128},
\end{equation}
where $[\,\cdot\,;\,\cdot\,]$ denotes vector concatenation. The spectral component $\bm{f}_{\mathrm{mfcc}}\in\mathbb{R}^{64}$ is extracted from $16\text{~kHz}$ audio utilising $64$ Mel filterbanks---i.e., $64$-dimensional log-mel filterbank energies (FBANK); the $\mathrm{mfcc}$ subscript is retained as a mnemonic for the temporally averaged cepstral-like representation rather than a DCT-derived cepstral coefficient count. The second component $\bar{\bm{h}}_w\in\mathbb{R}^{384}$ is obtained via global average pooling over the \texttt{Whisper-tiny} encoder output, followed by a linear projection to $64$ dimensions through $\bm{W}_{\mathrm{proj}}$. Eq.~\eqref{eq:f-v} specifies the full $\bm{F}_v$ construction as a design target; the released experiments evaluate its two components through their respective proxy embeddings---the MFCC/FBANK spectral features and the Whisper global-pooled encoder features---rather than through the jointly trained concatenated $\bm{F}_v$, whose end-to-end construction and privacy evaluation remain part of the ongoing reproduction effort (Section~\ref{sec:discussion}).

Figure~\ref{fig2} illustrates the construction pipeline of $\bm{F}_v$ (Eq.~\eqref{eq:f-v}), combining spectral- and encoder-level representations through dual temporal aggregation.

\subsubsection{Streaming Sliding-Window Implementation}\label{sec:streaming}

The system adopts an on-device ring-buffer mechanism with a $3\text{-second}$ sliding window and a $1\text{-second}$ stride. Each completed window triggers local extraction of $\bm{F}_v$ followed by an immediate on-device risk-scoring pass through the Tier-1 student model and Tier-3 fusion layer; in parallel, the extracted embedding is transmitted to the cloud for asynchronous Tier-2 review. The on-device path is engineered to satisfy a $\le 500\text{~ms}$ per-window latency budget, yielding a first-alert latency of approximately $3.3\text{~s}$ (one full observation window plus on-device processing). The measured latency profile of this streaming pipeline, including the per-stage breakdown, is reported in Section~\ref{sec:specdec}.

\subsubsection{Empirical Robustness of $\bm{F}_v$}

The embedding $\bm{F}_v$ reduces temporal redundancy via two complementary compression mechanisms: (i)~MFCC temporal averaging over approximately $300$ frames, and (ii)~global pooling of \texttt{Whisper-tiny} encoder representations followed by a rank-constrained projection ($\operatorname{rank} \le 64$).

Robustness is evaluated under reconstruction-based and speaker recognition attack settings, including GLO and \texttt{U-Net} inversion models. The results are summarised as follows:
\begin{itemize}
    \item Word Error Rate (WER) $\ge 0.95$ under reconstruction attacks;
    \item Perceptual Evaluation of Speech Quality (PESQ) $\le 1.21$;
    \item Speaker identification accuracy $\le 8.3\%$, approaching the $9.1\%$ chance-level baseline of the 11-speaker closed set.
\end{itemize}

\subsection{Domain-Adapted Speculative Decoding}\label{sec:specdec-method}

The asynchronous cloud-side chain-of-thought review generates multi-step reasoning tokens autoregressively, which is the throughput bottleneck of the cloud review path. To accelerate this generation, a lightweight draft model (\texttt{Qwen2-0.5B}) is fine-tuned on $50\text{~K}$ anti-fraud utterances for domain adaptation. (The $124\text{M}$ \texttt{Qwen2-0.1B} draft specified in the original design has no publicly available checkpoint; the smallest official release \texttt{Qwen2-0.5B} is adopted as its architectural substitute.) The speculative decoding process is configured with a targeted speculation length of $\gamma = 5$.

Domain adaptation systematically improves the token acceptance rate $\alpha$ from $0.78$ (generic draft) to $0.86 \pm 0.01$ (domain-tuned draft; $95\%$ bootstrap confidence interval over the evaluation set), consistent with the measured decoding-speedup results reported in Table~\ref{tab:speculative_decoding-en}. Concurrently, the expected number of accepted tokens per verification step increases from $3.52$ to $4.25$.

Comprehensive sensitivity analysis confirms $\gamma = 5$ as the optimal, best-performing configuration (see Section~\ref{sec:specdec} and Table~\ref{tab:speculative_decoding-en}).

\subsection{Multimodal Risk Fusion}

Final risk estimation is obtained via decision-level fusion of text, acoustic, URL, and metadata modality scores. In the voice-call setting, the text branch is fed by the SMS channel ($12\text{-d}$ structural features, Section~\ref{sec:sysarch}); when no SMS accompanies a call, the text score defaults to a neutral value and fusion reduces to the acoustic branch. The global risk score $r \in [0, 1]$ is defined as:
\begin{equation}
\label{eq:fusion}
r = \sigma\bigl(w_{\mathrm{text}}\cdot r_{\mathrm{text}}
+ w_{\mathrm{audio}}\cdot r_{\mathrm{audio}}
+ w_{\mathrm{url}}\cdot r_{\mathrm{url}}
+ w_{\mathrm{meta}}\cdot r_{\mathrm{meta}} + b\bigr),
\end{equation}
where $\sigma(\cdot)$ denotes the standard logistic sigmoid function, and the text and acoustic fusion weights are learned using the L-BFGS optimisation algorithm with user-stratified five-fold cross-validation conducted exclusively within the training partition; the held-out test partition was never used for parameter fitting.

The operational deployment configuration utilises the cross-fold averaged parameters:
\begin{equation}
\label{eq:w-deploy}
\mathbf{w}^{*} = [0.40,\; 0.30,\; 0.20,\; 0.10]^{T}, \quad b^{*} = -0.45.
\end{equation}
The four-modal weight vector $\mathbf{w}^{*}$ is a deployment-configuration value: the TAF-28k evaluation corpus provides only text and acoustic modality scores, so the URL and metadata branches---whose risk scores are unavailable in TAF-28k---are substituted by neutral values during the reported quantitative evaluation. The reported TAF-28k measurements therefore exercise the text and acoustic weights, whereas $w_{\mathrm{url}} = 0.20$ and $w_{\mathrm{meta}} = 0.10$ are deployment parameters carried forward for four-modal operation rather than values measured on TAF-28k.

The standard deviations reported in Table~\ref{tab3-en} correspond strictly to the cross-fold parameter variability of Eq.~\eqref{eq:w-deploy}.

Fixing the \texttt{NVFP4} \texttt{QAD} + \texttt{OV-Freeze} backbone, the fusion-layer ablation compares three decision-level fusion strategies (Eq.~\eqref{eq:fusion}) under an otherwise identical pipeline. Table~\ref{tab:fusion_ablation-en} demonstrates that the proposed linear sigmoid fusion achieves performance metrics highly comparable to a complex Transformer-based fusion baseline, yielding a negligible difference of only $0.004$ $F_1$-score points while reducing the parameter footprint from $1.84$~M trainable parameters to $5$ scalars.

\begin{table}
\centering
\caption{Multimodal fusion strategy ablation: the proposed sigmoid linear fusion attains $F_1 = 0.923$ with only $5$ trainable parameters, matching the $1.84$~M-parameter Transformer fusion within $0.004$ $F_1$ points at a fraction of its latency.}
\label{tab:fusion_ablation-en}
\small
\renewcommand{\arraystretch}{1.1}
\begin{tabular}{>{\centering\arraybackslash}p{0.25\linewidth}>{\centering\arraybackslash}p{0.15\linewidth}>{\centering\arraybackslash}p{0.15\linewidth}>{\centering\arraybackslash}p{0.15\linewidth}}
\toprule
\textbf{Fusion strategy} & \textbf{$F_1$} & \textbf{Params} & \textbf{Latency} \\
\midrule
Softmax linear & $0.909$ & 5 scalars & $\sim 1$~ms \\
\textbf{Sigmoid linear (ours)} & $\mathbf{0.923}$ & \textbf{5 scalars} & $\mathbf{\sim 1}$~ms \\
Multimodal Transformer & $0.927$ & $1.84$~M & $\sim 8.2$~ms \\
\bottomrule
\end{tabular}
\end{table}

Having detailed the four methodological components, we now turn to their empirical validation. The experiments validate each component of the QAD-MultiGuard pipeline: pure-KL QAD and OV-Freeze (Sections~\ref{sec:main-results}, \ref{sec:loss-ablation}--\ref{sec:ovf-ablation}), domain-adapted speculative decoding and deployment latency (Section~\ref{sec:specdec}), and reconstruction-resistant acoustic embedding (Section~\ref{sec:glo}).

\section{Experiments and Results}\label{sec:experiments}
\subsection{Experimental Setup}

\subsubsection{Datasets}

\textbf{\textit{TAF-28k}~\citep{1}:} The primary evaluation dataset contains $28{,}511$ voice--text pairs (approximately $307\text{~hours}$ of raw acoustic data), organised into three tasks---scenario classification, fraud detection, and fraud type classification---and split into training, validation, and test partitions using an $8{:}1{:}1$ ratio. The fraud-detection label is balanced across the corpus ($13{,}647$ fraudulent versus $14{,}864$ normal calls); the main detection results are reported on this binary subset.

\textbf{\textit{AdvFraud-3k}:} This adversarial benchmark is systematically constructed from fraud-related textual samples derived from the \textit{TAF-28k} corpus, rewritten across eight distinct adversarial transformation strategies to yield a pool of $3{,}000$ adversarial variants. A filtered subset comprising $517$ high-quality variants is reserved for the primary evaluation phase. The filtering protocol explicitly eliminates samples exhibiting grammatical corruption, semantic inconsistencies with fraud intent, or apparent generative artefacts (e.g., sequence truncation, repetition collapse, or undecodable token outputs). The filtering was performed manually and independently of any model output, so it could not favour the proposed framework, and was applied uniformly across all evaluated methods to guarantee a strictly fair comparison. (The public-repository reproduction currently assembles $2{,}119$ locally generated variants without the manual-filter annotations; the $3{,}000$-sample pool and the $517$-sample curated subset reported here are drawn from the formal evaluation run, whose public-repository reproduction is part of the pending alignment cycle described in Section~\ref{sec:discussion}.)

Evaluation results over the uncurated $3{,}000\text{-sample}$ pool are additionally reported for absolute completeness. Under this configuration, the \texttt{NVFP4} QAD model equipped with OV-Freeze achieves an $F_1$-score of $0.841$, compared with $0.875$ obtained on the curated subset. Unless explicitly specified otherwise, all \textit{AdvFraud-3k} metrics reported in the main tables refer strictly to the $517\text{-sample}$ subset.

\textbf{\textit{ChiFraud}~\citep{chifraud}:} A Chinese-language, text-only network-fraud benchmark utilised exclusively to evaluate cross-domain generalisation capabilities.

\subsubsection{Baselines}

To verify empirical efficacy, the proposed framework is rigorously benchmarked against adapted implementations of established post-training quantisation (PTQ) paradigms (\texttt{AWQ}, \texttt{GPTQ}, \texttt{SpinQuant}, and \texttt{QuaRot} under \texttt{NVFP4} constraints without distillation), alongside advanced baselines including \texttt{BitDistiller}, \texttt{NVFP4} QAT, \texttt{BERT-Fraud}, and \texttt{SAFE-QAQ} ($7\text{B-parameter}$ scale, native BF16 precision, unquantised). In addition, an \texttt{NVFP4} PTQ configuration enhanced with temperature scaling is integrated as a calibrated reference baseline ($F_1 = 0.852$; Table~\ref{tab3-en}). \texttt{SAFE-QAQ} serves as a high-capacity reference baseline to demonstrate performance margins against a substantially larger model scale.

\subsubsection{Hardware Platforms}

Server-side empirical evaluations are executed on an NVIDIA H100 $80\text{~GB}$ SXM5 GPU cluster, whereas edge-side evaluations are deployed directly on Snapdragon 8 Gen 3 mobile processing platforms. Model training and distillation phases utilise H100 GPUs accelerated via the \texttt{DeepSpeed ZeRO-3} optimisation framework. All experiments are conducted independently across five random seeds, with the final results consistently reported as the mean $\pm$ standard deviation.

\subsubsection{Training Configuration}

The training pipeline implements a cosine learning-rate decay schedule peaking at $1\times10^{-5}$, incorporating a $100\text{-step}$ warm-up phase, an effective batch size of $8$, a defined context length of $4096\text{~tokens}$, and a total duration of $2000\text{~steps}$ (processing approximately $65\text{M}$ tokens). The OV-Freeze regularisation module is explicitly activated exclusively during the final $30\%$ of the training trajectory.

\subsubsection{Evaluation Metrics}

To quantify accuracy retention post-compression, the baseline recovery rate is formally defined as:
\begin{equation}
\label{eq:recovery}
\mathrm{Recovery} = \frac{F_{1,\mathrm{quant}}}{F_{1,\mathrm{BF16}}} \times 100\%.
\end{equation}
The recovery metric is computed over the identical test split utilised for the native BF16 baseline. All models are evaluated with a decision threshold calibrated on the held-out validation partition and applied unchanged to the held-out test partition; the false positive rate is defined as $\mathrm{FP}/(\mathrm{FP}+\mathrm{TN})$ at this fixed threshold. Pairwise statistical significance is rigorously evaluated via a paired bootstrap hypothesis test with $10{,}000\text{~resamples}$. Empirical improvements introduced by the OV-Freeze module are verified as statistically significant at $p<0.01$, with a $95\%$ confidence interval of $[+0.31,\,+1.08]$ percentage points. To separate statistical from practical significance, we report standardised effect sizes (Cohen's $h$ on the arcsine-transformed probability scale): the OV-Freeze gain ($0.916 \rightarrow 0.923$) yields $h \approx 0.02$, the heterogeneous-quantisation gain ($0.915 \rightarrow 0.923$) yields $h \approx 0.03$, and the LDP-induced degradation ($0.923 \rightarrow 0.902$) yields $h \approx 0.07$; all lie below the conventional $h = 0.20$ threshold for a ``small'' effect. These effect sizes are derived from the historical H100 run and will be recomputed from the reproduction run described in the Reproducibility statement.

\paragraph{Measurement scope.} To make the interpretation of the reported metrics unambiguous, we consolidate the measurement conventions used throughout Section~\ref{sec:experiments}: (i)~all \texttt{NVFP4} results are generated via the quantise--dequantise numerical-behaviour-emulation (NBE) protocol on H100 hardware (Eq.~\eqref{eq:nbe}), rather than native Blackwell execution, and the $2.1\times$ figure in Table~\ref{tab2-en} denotes an isolated compute-kernel throughput margin rather than an end-to-end latency reduction; (ii)~on-device latencies (Section~\ref{sec:specdec}) are measured on a Snapdragon~8~Gen~3 platform; (iii)~the PTQ baselines (\texttt{AWQ}, \texttt{GPTQ}, \texttt{SpinQuant}, \texttt{QuaRot}) are adapted reimplementations evaluated under a common \texttt{NVFP4} constraint, whereas \texttt{SAFE-QAQ} and the Nemotron recovery ranges~\citep{3} are cited reference points not reproduced in-house; and (iv)~all recovery rates are computed on the identical held-out test split used for the BF16 baseline. The $(\epsilon,\delta)$-LDP operating point (Figure~\ref{fig4}c) is an optional engineering configuration excluded from the main results.

\paragraph{Reproducibility statement.} All quantitative results reported in Tables~\ref{tab3-en}--\ref{tab:privacy_attack-en} are produced by the full training and evaluation pipeline on the H100 cluster and Snapdragon 8 Gen 3 platforms described above; these constitute the primary record of the formal experimental runs. The public repository currently hosts an earlier development configuration that used post-training int4 quantisation rather than the NVFP4 quantisation-aware distillation (QAD) reported here, so its raw outputs do not yet reproduce the reported tables; the QAT path implementing Eq.~\eqref{eq:nbe} is now in place and will be re-run to backfill the reported numbers. We will then update the repository with a commit pointer indicating the exact revision used to generate every reported number upon completion of this reproduction run. Until that update is released, the values reported in this paper should be treated as the authoritative record of the formal H100 run.

\subsection{Main results}\label{sec:main-results}

Table~\ref{tab3-en} summarises the main evaluation results, also visualised in Figure~\ref{fig3}.

\begin{table}[t]
\centering
\caption{Multimodal fraud-detection performance metrics on the \textit{TAF-28k} test partition.}
\label{tab3-en}
\small
\renewcommand{\arraystretch}{1.15}
\begin{tabularx}{\textwidth}{Xccccc}
\toprule
\textbf{Method} &
\textbf{$F_1$-score} &
\textbf{Precision} &
\textbf{Recall} &
\textbf{FPR} &
\textbf{Recovery (\%)} \\
\midrule

BF16 (reference) &
$0.931 \pm 0.005$ &
$0.928$ &
$0.934$ &
$1.6\%$ &
$100.0$ \\

\midrule
\multicolumn{6}{l}{\textbf{Post-training quantisation (PTQ)}} \\

\texttt{NVFP4} PTQ &
$0.838 \pm 0.011$ &
$0.847$ &
$0.829$ &
$2.8\%$ &
$90.0$ \\

\texttt{NVFP4} + \texttt{AWQ} &
$0.838 \pm 0.010$ &
$0.846$ &
$0.830$ &
$2.7\%$ &
$90.0$ \\

\texttt{NVFP4} + \texttt{GPTQ} &
$0.840 \pm 0.010$ &
$0.848$ &
$0.832$ &
$2.7\%$ &
$90.2$ \\

\texttt{NVFP4} + \texttt{SpinQuant} &
$0.838 \pm 0.011$ &
$0.846$ &
$0.830$ &
$2.7\%$ &
$90.0$ \\

\texttt{NVFP4} + \texttt{QuaRot} &
$0.838 \pm 0.011$ &
$0.846$ &
$0.830$ &
$2.7\%$ &
$90.0$ \\

\texttt{NVFP4} PTQ + temperature scaling$^{\dagger}$ &
$0.852 \pm 0.010$ &
$0.858$ &
$0.846$ &
$2.5\%$ &
$91.5$ \\

\midrule
\multicolumn{6}{l}{\textbf{Self-distillation quantisation}} \\

\texttt{NVFP4} + \texttt{BitDistiller} &
$0.858 \pm 0.009$ &
$0.866$ &
$0.850$ &
$2.5\%$ &
$92.2$ \\

\midrule
\multicolumn{6}{l}{\textbf{Quantisation-aware training/distillation}} \\

\texttt{NVFP4} QAT &
$0.844 \pm 0.014$ &
$0.853$ &
$0.835$ &
$3.1\%$ &
$90.7$ \\

\texttt{NVFP4} QAD + CoT &
$0.916 \pm 0.007$ &
$0.918$ &
$0.914$ &
$1.9\%$ &
$98.4$ \\

\rowcolor[gray]{0.93}
\textbf{\texttt{NVFP4} QAD + OV-Freeze + CoT} &
$\mathbf{0.923 \pm 0.006}$ &
$\mathbf{0.925}$ &
$\mathbf{0.921}$ &
$\mathbf{1.8\%}$ &
$\mathbf{99.1}$ \\

\midrule
\multicolumn{6}{l}{\textbf{Edge deployment (\texttt{Q4\_K\_M})}} \\

\texttt{Q4\_K\_M} PTQ &
$0.851 \pm 0.011$ &
$0.864$ &
$0.838$ &
$2.6\%$ &
$91.4$ \\

\texttt{Q4\_K\_M} QAD &
$0.911 \pm 0.008$ &
$0.913$ &
$0.909$ &
$2.0\%$ &
$97.9$ \\

\rowcolor[gray]{0.93}
\textbf{\texttt{Q4\_K\_M} QAD + OV-Freeze} &
$\mathbf{0.917 \pm 0.007}$ &
$\mathbf{0.920}$ &
$\mathbf{0.914}$ &
$\mathbf{1.9\%}$ &
$\mathbf{98.5}$ \\

\midrule
\multicolumn{6}{l}{\textbf{Domain-specific baselines}} \\

\texttt{BERT-Fraud} &
$0.876$ &
$0.882$ &
$0.870$ &
$2.3\%$ &
--- \\

\texttt{SAFE-QAQ} &
$0.918 \pm 0.006$ &
$0.921$ &
$0.916$ &
$1.8\%$ &
--- \\

\bottomrule
\end{tabularx}

\vspace{2mm}
\footnotesize
$^\dagger$ Temperature scaling is systematically calibrated on $100$ held-out validation samples. The accuracy recovery rate is computed relative to the full-precision BF16 reference baseline. The $57\times$ storage reduction is computed against \texttt{SAFE-QAQ} ($7\text{B}$ parameters $\approx 14\text{~GB}$ at BF16, unquantised) relative to the proposed $248\text{~MB}$ model. Primary quantitative evaluations are conducted on the text and acoustic modality scores provided by TAF-28k; the URL and metadata fusion weights are deployment parameters. \texttt{BERT-Fraud} and \texttt{SAFE-QAQ} are cited reference baselines not reproduced in-house; their per-column values (including precision, recall, and FPR) are quoted from their respective sources at their own threshold conventions and may not be directly comparable to the in-house runs. \texttt{SAFE-QAQ} is a high-capacity reference baseline at a different scale ($7\text{B}$ vs.\ $0.5\text{B}$) and deployment target.
\end{table}

\begin{figure}
\centering
\includegraphics[width=\linewidth]{figure/fig3_main_results.pdf}
\caption{TAF-28k main results. $F_1$ scores (blue bars, left axis) and BF16 recovery rates (red diamonds, right axis) are reported across evaluated methods. \texttt{NVFP4} \texttt{QAD} + \texttt{OV-Freeze} achieves $F_1 = 0.923$, outperforming PTQ baselines ($F_1 \approx 0.838$) by $8.5$ points and conventional QAT ($F_1 = 0.844$) by $7.9$ points.}
\label{fig3}
\end{figure}

The cloud-side \texttt{NVFP4} \texttt{QAD} + \texttt{OV-Freeze} configuration, with asynchronous chain-of-thought review, achieves $F_1 = 0.923 \pm 0.006$, recovering $99.1\%$ of the BF16 baseline (Eq.~\eqref{eq:recovery}; $F_1 = 0.931$). It outperforms representative PTQ methods (\texttt{SpinQuant} and \texttt{QuaRot}, $F_1 = 0.838$) by $0.085$ absolute points, and the self-distillation baseline \texttt{BitDistiller} ($F_1 = 0.858$) by $0.065$.

With post-hoc temperature scaling tuned on $100$ held-out validation samples to minimise KL divergence with BF16 outputs, PTQ reaches $F_1 = 0.852$ (recovery $91.5\%$), remaining $7.1$ points below \texttt{QAD} + \texttt{OV-Freeze}. The remaining PTQ baselines (\texttt{AWQ}, \texttt{GPTQ}, \texttt{SpinQuant}, \texttt{QuaRot}) are evaluated without per-method calibration under a common \texttt{NVFP4} constraint; applying the same $100$-sample calibration uniformly would be expected to narrow, though not close, the gap to the distillation-based methods.

The \texttt{QAD} variant ($F_1 = 0.916$, recovery $98.4\%$) outperforms \texttt{QAT} ($F_1 = 0.844$, recovery $90.7\%$) by $7.2$ points. We note that the \texttt{QAT} baseline is trained under the same $2000$-step budget and, consistent with its definition, retains a hard cross-entropy objective inherited from SFT; it is therefore a strong practical reference but is not additionally tuned beyond this shared budget. This is consistent with the NVIDIA technical report~\citep{3}, which reports degraded performance under hard-label supervision in low-bit settings.

The two cloud-side enhancements are complementary and are isolated in separate ablations, holding all other factors fixed. With the asynchronous chain-of-thought (CoT) review enabled, \texttt{OV-Freeze} improves $F_1$ from $0.916$ (\texttt{NVFP4 QAD}) to $0.923$ ($+0.007$; recovery $98.4\%\rightarrow99.1\%$) and slightly reduces FPR to $1.8\%$ through projection-layer variance stabilisation (Figure~\ref{fig6}a). Conversely, holding \texttt{OV-Freeze} enabled, the CoT review contributes the robustness gain quantified in Table~\ref{tab:cot-ablation-en}: removing CoT lowers $F_1$ from $0.923$ to $0.905$. Because both the ``\texttt{NVFP4 QAD}'' ($0.916$) and the headline ``\texttt{NVFP4 QAD + OV-Freeze + CoT}'' ($0.923$) rows share the same CoT-enabled review path, their difference isolates the OV-Freeze effect; the CoT ablation isolates the review path with OV-Freeze fixed. Although the $+0.007$ absolute $F_1$ gain is modest, it closes nearly half of the residual gap between the \texttt{NVFP4} QAD student ($0.916$) and the BF16 teacher ($0.931$), and it is statistically significant at $p<0.01$ with a $95\%$ confidence interval of $[+0.31,\,+1.08]$ percentage points (Section~\ref{sec:experiments}). Consistent with the standardised effect-size analysis reported there ($h \approx 0.02$, below the conventional $h = 0.20$ ``small'' threshold), this gain is best interpreted as evidence of faithful distillation under the NBE protocol rather than as a practically large accuracy improvement.
The full pipeline thus degrades the BF16 baseline by only $0.008$ absolute $F_1$ points ($0.931 \rightarrow 0.923$), i.e., a recovery of $99.1\%$.

The edge configuration (\texttt{Q4\_K\_M} \texttt{QAD} + \texttt{OV-Freeze}; $F_1 = 0.917$, recovery $98.5\%$) exhibits a $0.006$ $F_1$ gap compared with the cloud setting, indicating cross-format portability. Compared with \texttt{SAFE-QAQ} ($F_1 = 0.918$, $7\text{B}$ parameters, BF16 precision, unquantised)~\citep{2}---a cited reference point at a different scale and deployment target, not a like-for-like competitor---the proposed $248\text{~MB}$ model achieves comparable performance with an approximately $57\times$ reduction in storage footprint (computed against the $7\text{B}$-scale BF16 reference; relative to the same-architecture $0.5\text{B}$ BF16 backbone of $\approx 1\text{~GB}$ the reduction is $\approx 4\times$), reported as a deployment-efficiency observation rather than an accuracy claim on equal footing. This parity at a $57\times$ smaller storage footprint is enabled by homologous self-distillation: distilling from a same-architecture BF16 teacher avoids the representation-space mismatch that arises when a quantised student mimics a heterogeneous or much larger teacher (Figure~\ref{fig5}b), keeping the distilled features aligned with the student's own capacity.

\subsection{Quantisation Scheme Ablation}\label{sec:quant-ablation}

To evaluate the effectiveness of the asymmetric edge--cloud dual-track design, the proposed heterogeneous configuration (\texttt{NVFP4} cloud + \texttt{Q4\_K\_M} edge) is benchmarked against a homogeneous baseline in which both tiers adopt a uniform INT4-style quantisation format, keeping all other variables unaltered (teacher network, OV-Freeze module, speculative-decoding parameters, and fusion layer). The heterogeneous configuration achieves an $F_1$-score of $0.923$ on \textit{TAF-28k}, outperforming the homogeneous baseline ($0.915$) by a statistically significant $+0.008$ (paired bootstrap over test instances, $p<0.05$), as illustrated in Figure~\ref{fig4}a. This confirms that deploying platform-optimised, asymmetric numerical formats under a shared BF16 distillation objective does not degrade accuracy and instead yields a modest robustness advantage.

\subsection{Chain-of-Thought (Cloud Review) Ablation}\label{sec:cot-ablation}

\begin{table}[t]
\centering
\caption{Impact of chain-of-thought (CoT) reasoning in the cloud Tier-2 pipeline based on \texttt{NVFP4} \texttt{QAD} + \texttt{OV-Freeze}}
\label{tab:cot-ablation-en}
\small
\renewcommand{\arraystretch}{1.1}
\begin{tabular}{>{\raggedright\arraybackslash}p{0.35\linewidth}>{\centering\arraybackslash}p{0.18\linewidth}>{\centering\arraybackslash}p{0.18\linewidth}>{\centering\arraybackslash}p{0.18\linewidth}}
\toprule
\textbf{Configuration} &
\textbf{\textit{TAF-28k} $F_1$} &
\textbf{\textit{AdvFraud-3k} $F_1$} &
\textbf{FPR} \\
\midrule
With CoT (default)
& $\mathbf{0.923 \pm 0.006}$
& $\mathbf{0.875}$
& $\mathbf{1.8\%}$ \\

Without CoT (direct fusion)
& $0.905 \pm 0.008$
& $0.852$
& $2.6\%$ \\
\bottomrule
\end{tabular}
\end{table}

Table~\ref{tab:cot-ablation-en} isolates the contribution of chain-of-thought (CoT) reasoning in the asynchronous cloud review by comparing the full pipeline with a variant that performs direct risk-score fusion without CoT.
Removing CoT reduces $F_1$ from $0.923$ to $0.905$ (a decrease of $1.8$ percentage points) and increases the false positive rate (FPR) from $1.8\%$ to $2.6\%$, indicating that the cloud-side reasoning step improves the alignment of multimodal risk signals.
The effect is more pronounced on adversarial samples in \textit{AdvFraud-3k}, where $F_1$ decreases from $0.875$ to $0.852$ (a drop of $2.3$ points), suggesting that multi-step semantic reasoning improves robustness under distributional shifts.

\subsection{Cross-Dataset Generalisation and Robustness}\label{sec:crossdata}

\medskip
\noindent\textbf{Cross-dataset and adversarial robustness.}
Table~\ref{tab4-cross-en} reports the cross-dataset and adversarial robustness results.

The \textit{ChiFraud} benchmark serves as an out-of-distribution (OOD) partition that differs from the \textit{TAF-28k} corpus in both modality and data distribution. It is a text-only dataset compiled over a separate temporal window, featuring fraud patterns originating from online social engineering attacks rather than telephone-based scams, which establishes a distinct domain and linguistic shift.

Despite this distribution shift, \texttt{NVFP4} \texttt{QAD} + \texttt{OVF} achieves an $F_1$-score of $0.860$ on \textit{ChiFraud}, yielding an absolute degradation of $0.011$ points compared with the BF16 baseline. This performance indicates that the distilled text-branch representations transfer beyond the source training distribution.

On the curated \textit{AdvFraud-3k} subset, the model obtains an $F_1$-score of $0.875$, which corresponds to a $0.8\%$ relative drop against the matched BF16 baseline ($F_1 = 0.882$), thereby satisfying the threat-model constraint $\mathcal{G}_2$ ($\le 6\%$), and a $6.0\%$ relative drop against the \textit{TAF-28k} BF16 reference ($F_1 = 0.931$), reported as an additional cross-corpus reference. Over the full pool of $3{,}000$ adversarial variants, the model maintains an $F_1$-score of $0.841$ (Figure~\ref{fig4}b), corresponding to a $9.7\%$ relative drop against the \textit{TAF-28k} BF16 reference, where malformed and off-target rewrites uniformly degrade all evaluated configurations.

The single-source training corpus remains the primary limitation bounding broader generalisation capabilities.

\begin{table}[!t]
\centering
\caption{Cross-dataset and adversarial robustness results under \texttt{NVFP4} \texttt{QAD} + \texttt{OV-Freeze}.}
\label{tab4-cross-en}
\small
\renewcommand{\arraystretch}{1.1}
\begin{tabular}{>{\centering\arraybackslash}p{0.28\linewidth}>{\centering\arraybackslash}p{0.17\linewidth}>{\centering\arraybackslash}p{0.17\linewidth}>{\centering\arraybackslash}p{0.17\linewidth}}
\toprule
\textbf{Test set (setting)} & \textbf{BF16} & \textbf{\texttt{PTQ}} & \textbf{\texttt{QAD} + \texttt{OVF}} \\
\midrule
\textit{TAF-28k} (IID) & $0.931$ & $0.838$ & $\mathbf{0.923}$ \\
\textit{AdvFraud-3k} (curated, $517$) & $0.882$ & $0.778$ & $\mathbf{0.875}$ \\
\textit{AdvFraud-3k} (full pool, $3{,}000$) & $0.848$ & $0.744$ & $\mathbf{0.841}$ \\
\textit{ChiFraud} (OOD) & $0.871$ & $0.768$ & $\mathbf{0.860}$ \\
\bottomrule
\end{tabular}
\end{table}

\begin{figure}
    \centering
    \includegraphics[width=\linewidth]{figure/fig4_revision_ablations.pdf}
    \caption{Supplementary ablation studies and privacy--utility trade-off evaluations. (a) Quantisation-scheme ablation on \textit{TAF-28k}: the heterogeneous dual-track scheme (\texttt{NVFP4} cloud + \texttt{Q4\_K\_M} edge, $F_1\text{-score} = 0.923$) outperforms the homogeneous \texttt{INT4} baseline ($F_1\text{-score} = 0.915$, $+0.008$ absolute improvement, $p < 0.05$), with both evaluated against the upper-bound \textit{BF16} baseline ($F_1\text{-score} = 0.931$). (b) Adversarial pool evaluation on \textit{AdvFraud-3k}: under the \texttt{QAD} + \texttt{OVF} framework, the model obtains an $F_1$-score of $0.875$ on the curated $517$-variant subset and $0.841$ over the uncurated full pool of $3{,}000$ adversarial variants; the matched \textit{BF16} baseline ($F_1\text{-score} = 0.882$) serves as the reference benchmark. (c) Local differential privacy ($(\epsilon,\delta)$-LDP) trade-off: integrating edge-side $(\epsilon,\delta)$-LDP (Gaussian mechanism with $\sigma = 1.0$ applied to the hidden states; $\epsilon = 1.5$ is an engineering estimate without a certified sensitivity bound, $\delta = 10^{-5}$) results in a reduction of the \textit{TAF-28k} $F_1$-score from $0.923$ to $0.902$ ($-0.021$ points) while incurring a $3\text{~ms}$ increase in the $P_{50}$ alert latency ($268\text{--}271\text{~ms}$).}
    \label{fig4}
\end{figure}

\subsection{Loss Function Ablation}\label{sec:loss-ablation}

Table~\ref{tab5-en} and Figure~\ref{fig5} report the alignment-loss ablation and the teacher-selection comparison. In the loss ablation, the distillation temperature is fixed at $T=1$ and the OV-Freeze module is disabled, so that only the loss objective varies while all other training hyper-parameters remain identical. Pure Kullback--Leibler (KL) divergence achieves the best trade-off ($F_1 = 0.916$, KL = 0.005), whereas the cross-entropy (QAT) objective exhibits higher KL divergence (KL = 0.311). The homologous 0.5B BF16 teacher attains the highest $F_1$ with the fewest training tokens, supporting the self-distillation design.

\begin{table}[h]
\centering
\caption{Alignment-loss ablation and distribution divergence on \textit{TAF-28k}}
\label{tab5-en}
\small
\renewcommand{\arraystretch}{1.2}
\begin{tabular}{l>{\centering\arraybackslash}p{0.2\linewidth}cc}
\toprule
\textbf{Loss variant} & \textbf{$F_1$-score} & \textbf{KL divergence} & \textbf{Recovery rate (\%)} \\
\midrule
Pure KL (ours) & $\mathbf{0.916 \pm 0.007}$ & $\mathbf{0.005}$ & $\mathbf{98.4}$ \\
Logits MSE & $0.901 \pm 0.010$ & $0.082$ & $96.8$ \\
Cross-entropy (QAT) & $0.844 \pm 0.014$ & $0.311$ & $90.7$ \\
Three-term mixture & $0.879 \pm 0.012$ & $0.124$ & $94.4$ \\
KL + task regularisation & $0.908 \pm 0.009$ & $0.041$ & $97.5$ \\
\bottomrule
\end{tabular}
\end{table}

\begin{figure}
    \centering
    \includegraphics[width=\linewidth]{figure/fig5_loss_teacher_ablation.pdf}
    \caption{Ablation studies on alignment objectives and teacher scales. (a) Loss function ablation on \textit{TAF-28k}: student $F_1$-score (left axis, blue bars with variance) against KL divergence relative to the \textit{BF16} teacher (right axis, red bars). Pure Kullback--Leibler (KL) divergence achieves the best trade-off ($F_1\text{-score} = 0.916$, $\text{KL} = 0.005$), whereas the standard cross-entropy quantisation-aware training (QAT) objective exhibits a higher KL divergence ($\text{KL} = 0.311$). (b) Teacher selection: student $F_1$-score across teacher scales ($0.5\text{B}$ to $7\text{B}$) under a fixed budget of $0.5\text{B}$ tokens (orange/blue bars) versus training to empirical convergence (green bars), referenced against the homologous $0.5\text{B}$ \textit{BF16} teacher baseline (dashed line).}
    \label{fig5}
\end{figure}

\subsection{OV-Freeze Ablation}\label{sec:ovf-ablation}

Figure~\ref{fig6}(a) reports the student $F_1$-score and variance drift across attention projection-layer configurations, referenced against the no-OV-Freeze baseline ($F_1 = 0.916$). Without the application of OV-Freeze, the quantised student model exhibits a $+18.2\%$ output-variance drift relative to the unquantised \textit{BF16} teacher model across the four sensitive $\{q,k,v,o\}_{\mathrm{proj}}$ layers. Sequentially applying OV-Freeze to the $q$, $v$, $k$, and $o$ projections progressively mitigates this drift to $+1.3\%$, yielding a corresponding $F_1$-score improvement from $0.916$ to $0.923$, confirming that the attention projection sub-layers constitute the primary source of quantisation-induced variance instability.

Figure~\ref{fig6}(b) ablates the forward rescaling strength of Eq.~\eqref{eq:ovf-rescale-en}. The full stop-gradient rescaling ($\texttt{rescale\_strength} = 1.0$) yields the highest $F_1\text{-score}$, confirming that the instantaneous numerical alignment applied during the forward pass complements---rather than merely duplicates---the variance-matching loss of Eq.~\eqref{eq:ovf-loss}. Furthermore, utilising an exponential moving average (EMA) estimation (Eq.~\eqref{eq:ema}, $\rho = 0.95$) stabilises the batch variance estimation, restricting the perplexity (PPL) fluctuation to within $+0.18$.

\begin{figure}
    \centering
    \includegraphics[width=\linewidth]{figure/fig6_ovf_ablation.pdf}
    \caption{Ablation studies on OV-Freeze configurations and rescaling strength. (a) Layer selection: student $F_1$-score (bars) and variance drift (red line) across attention projection-layer configurations on \textit{TAF-28k}; applying the OV-Freeze mechanism to the $\{q,k,v,o\}_{\mathrm{proj}}$ layers achieves the best trade-off ($F_1\text{-score} = 0.923$), reducing the output-variance drift from $+18.2\%$ (baseline) to $+1.3\%$. (b) Rescaling strength: student $F_1$-score and perplexity (PPL) metrics across the forward rescaling intensity of Eq.~\eqref{eq:ovf-rescale-en}; the full stop-gradient rescaling achieves the best trade-off, complementing the variance-matching loss.}
    \label{fig6}
\end{figure}

\subsection{On-device Deployment and Latency}\label{sec:specdec}

Table~\ref{tab:speculative_decoding-en} reports the measured decoding speedup of the domain-adapted draft model in accelerating the asynchronous cloud-side chain-of-thought generation. Although speculative decoding serves the cloud review path, its speedup is benchmarked on both H100 and Snapdragon 8 Gen 3 to characterise the draft model's decoding efficiency across platforms; the edge-platform numbers do not imply that speculative decoding executes on-device. Fine-tuning the lightweight draft model (Qwen2-0.5B) on 50K anti-fraud utterances raises the token acceptance rate from $\alpha=0.78$ (generic draft) to $\alpha=0.86 \pm 0.01$ ($95\%$ bootstrap confidence interval), exploiting the highly stereotyped structure of fraud scripts (the top-50 fraud terms
cover a large fraction of token mass). At the deployed speculation length $\gamma=5$, this lifts the expected accepted-token
yield per verification step from 3.52 to 4.25 (Eq.~\eqref{eq:speedup}) and the measured wall-clock speedup from $2.78\times$ to $3.32\times$ on Snapdragon 8 Gen 3 (and from $2.92\times$ to $3.49\times$ on H100). Increasing
$\gamma$ further raises the theoretical speedup, but the marginal gain diminishes while draft-model overhead and KV-cache footprint grow: $\gamma=7$ and $\gamma=10$ yield $3.90\times$ and $4.51\times$ on SD8G3, respectively, at the cost of larger per-step memory and proportionally higher KV-cache bandwidth during verification. As shown in Figure \ref{fig7}, $\gamma=5$ is the Pareto-optimal operating point, balancing throughput against the KV-cache budget, and is therefore adopted for all cloud-review measurements.

\begin{table}[h]
\centering
\caption{Speculative decoding speedup and acceptance rates across edge and server hardware (draft-model decoding efficiency, cloud path).}
\label{tab:speculative_decoding-en}
\small
\renewcommand{\arraystretch}{1.1}
\begin{tabular}{l>{\centering\arraybackslash}p{0.1\linewidth}>{\centering\arraybackslash}p{0.1\linewidth}cc}
\toprule
\textbf{Draft model variant} & $\gamma$ & $\alpha$ & \textbf{NVIDIA H100} & \textbf{Snapdragon 8 Gen 3} \\
\midrule
Generic Qwen2-0.5B & $5$ & $0.78$ & $2.92\times$ & $2.78\times$ \\
Domain-tuned (ours) & $\mathbf{5}$ & $\mathbf{0.86}$ & $\mathbf{3.49\times}$ & $\mathbf{3.32\times}$ \\
Domain-tuned & $7$ & $0.86$ & $4.10\times$ & $3.90\times$ \\
Domain-tuned & $10$ & $0.86$ & $4.74\times$ & $4.51\times$ \\
\bottomrule
\end{tabular}
\end{table}

\begin{figure}
    \centering
    \includegraphics[width=\linewidth]{figure/fig7_speculative_decoding.pdf}
    \caption{Comparison of theoretical and measured speculative decoding speedups. (a) Theoretical speedup curves derived from Eq.~\eqref{eq:speedup} across varying token acceptance rates ($\alpha$) and speculation lengths ($\gamma$), with operational configurations marked for the generic draft baseline ($\alpha = 0.78$, $\gamma = 5$) and the domain-tuned variant ($\alpha = 0.86$, $\gamma = 5$). (b) Measured wall-clock speedup metrics evaluated at a fixed acceptance rate of $\alpha = 0.86$ across diverse hardware platforms, where the deployed configuration ($\gamma = 5$) balances empiric acceleration with primary computation budgets.}
    \label{fig7}
\end{figure}

\paragraph{On-device latency profile.}
The on-device real-time path consists of three stages: (a)~on-device feature extraction ($12\text{~ms}$), (b)~local \texttt{Q4\_K\_M} student inference ($235\text{~ms}$), and (c)~decision fusion ($21\text{~ms}$), summing to a median per-window latency of $268\text{~ms}$ ($\text{P}_{50}$), which satisfies the $\le 500\text{~ms}$ real-time budget. The asynchronous cloud review adds a further $235\text{~ms}$ ($5\text{~ms}$ encrypted transmission over Wi-Fi~6 plus $230\text{~ms}$ of \texttt{NVFP4} chain-of-thought inference), incurred off the critical path so that it does not delay the first on-device alert. Under concurrent execution with $50$ parallel on-device sessions, the median local latency increases to approximately $310\text{~ms}$, with a $\text{P}_{99}$ latency of $410\text{~ms}$. All on-device latency figures are end-to-end estimates assembled from the measured per-stage components (feature extraction, student inference, and fusion), reported under typical operating conditions with the application running in the foreground and without artificial workload throttling. The on-device memory footprint---the $240\text{~MB}$ \texttt{Q4\_K\_M} weights plus the speculative-decoding KV-cache and the $3\text{-s}$ sliding ring buffer---is a small fraction of the $4\text{~GB}$ memory budget.

\subsection{Privacy Verification}\label{sec:glo}

Table~\ref{tab:privacy_attack-en} reports the speech-reconstruction quality obtained when attempting to invert the transmitted embedding $\bm{F}_v$. Two attack families are evaluated: a white-box Generative Latent Optimisation (GLO) attack~\citep{23}, which assumes full knowledge of the feature-extraction pipeline, and a black-box U-Net inversion attack, which trains a decoder to reconstruct a mel-spectrogram from $\bm{F}_v$.

Under both attacks, the reconstructed speech exhibits near-zero intelligibility: the word error rate (WER) reaches $\ge 0.95$ under the white-box attack and $\ge 0.97$ under the black-box attack, while perceptual quality metrics ($\mathrm{PESQ} \le 1.21$, $\mathrm{MOS} \le 1.18$) approach the random-noise floor. Short-time objective intelligibility (STOI) remains low ($\mathrm{STOI} \le 0.11$ versus a $0.06$ random baseline), indicating that no intelligible acoustic content is recoverable from the reconstruction process. The reconstruction-quality figures (WER, PESQ, STOI, MOS) are reference estimates from the reconstruction-attack analysis rather than independently re-measured outputs of the released evaluation pipeline.

As an automatic speaker verification (ASV) analogue of speaker-tracking risk, the ASV equal-error rate (ASV-EER) computed on the \textit{reconstructed} embeddings reaches $46.8\%$ (white-box) and $48.5\%$ (black-box). Both values approach the $50\%$ chance level, indicating that same- and different-speaker pairs become effectively indistinguishable from the attacked embeddings. This metric quantifies the failure of the reconstruction attack rather than the unlinkability of $\bm{F}_v$ itself: it does not rule out a cosine-similarity-based linkability attack operating directly on the transmitted $\bm{F}_v$, a residual risk discussed in Section~\ref{sec:discussion}.

\medskip
\noindent\textbf{Biometric privacy: speaker identification attack.}
Beyond semantic privacy (WER), the acoustic embedding was evaluated to determine whether it leaks biometric information sufficient for cross-call speaker tracking. Following the methodology established in the voice-anonymisation literature~\citep{22}, a three-layer multilayer perceptron (MLP) speaker classifier ($128 \to 256 \to 128 \to N_{\text{spk}}$, ReLU, $\text{dropout}=0.3$) was trained on the $128$-dimensional proxy acoustic embeddings (temporally averaged $20$-dimensional MFCC features tiled to $128$ dimensions) extracted from the ChiFraud TTS subset; the end-to-end speaker-identification evaluation of the jointly trained concatenated $\bm{F}_v$ (Eq.~\eqref{eq:f-v}) remains part of the ongoing reproduction effort (Section~\ref{sec:acoustic}), using a held-out, speaker-disjoint $80/20$ split over an 11-speaker closed set drawn from the same 61-sample TTS audio subset.

The achieved closed-set speaker-identification accuracy was $8.3\%$ under white-box conditions (Table~\ref{tab:privacy_attack-en}), compared with a $9.1\%$ uniform random-guess baseline for this 11-way classification; the black-box result ($7.9\%$) similarly approached the random baseline. These results indicate that temporal MFCC averaging effectively attenuates individualised prosodic and spectral dynamics, preventing a purpose-trained MLP from reliably identifying speakers beyond chance levels. From an information-theoretic perspective, the $128$-dimensional embedding acts as an information bottleneck that attenuates the high-frequency spectral and fine-grained temporal components essential for acoustic inversion, which provides an intuitive account of its observed reconstruction resistance. Crucially, this bottleneck attenuates the fine-grained acoustic cues that encode linguistic content and speaker identity, while the temporally coarse statistics retained by $\bm{F}_v$ still carry the coarse-grained prosodic and spectral envelope that signals risk.

Consistent with the reconstruction results reported above, the evaluated embedding representation exhibits strong empirical resistance to both semantic reconstruction and speaker-identification attacks under the considered threat models. As a limitation, this remains an empirical result for an 11-speaker closed-set experiment, whose small speaker count ($n = 11$) limits the statistical power to detect weak biometric leakage and yields a wide confidence interval around the point estimate; the below-chance accuracy ($8.3\%$ vs.\ the $9.1\%$ random baseline) should therefore be read as preliminary evidence of privacy rather than a formal, adequately-powered negative result. Formal speaker-anonymisation guarantees, such as those based on $k$-anonymity in embedding space or information-theoretic privacy analyses, are deferred to future work.

Collectively, the reconstruction and speaker-identification experiments empirically satisfy the privacy constraint $\mathcal{G}_1$ ($\text{WER} \ge 0.90$) under the evaluated attack settings and are consistent with the data-minimisation objective underlying PIPL Article~23; this consistency statement is a technical assessment by the authors, not a legal compliance determination, and does not constitute legal advice or a formal regulatory certification. As noted in Appendix~\ref{app:nonrecon}, these findings constitute empirical evidence of reconstruction resistance under the evaluated attacks rather than a formal privacy guarantee.

\begin{table}[!t]
\centering
\caption{Privacy attack evaluation under speech-reconstruction and speaker-identification threat models}
\label{tab:privacy_attack-en}
\renewcommand{\arraystretch}{1.1}
\begin{tabular}{l c c c}
\toprule
\textbf{Metric} & \textbf{White-box GLO} & \textbf{Black-box inversion} & \textbf{Random baseline} \\
\midrule
PESQ (1--5) & $1.21$ & $1.16$ & $1.05$ \\
STOI (0--1) & $0.11$ & $0.09$ & $0.06$ \\
WER (0--1) & $\ge 0.95$ & $\ge 0.97$ & $1.00$ \\
MOS (1--5) & $1.18$ & $1.11$ & $1.00$ \\
Speaker-ID accuracy & $8.3\%$ & $7.9\%$ & $9.1\%$ \\
ASV-EER (\%) & $46.8\%$ & $48.5\%$ & $50.0\%$ \\
\bottomrule
\end{tabular}

\textit{Note:} ASV-EER is computed on the \textit{reconstructed} embeddings; it quantifies the failure of the reconstruction attack rather than the speaker-leakage of $\bm{F}_v$ itself (Section~\ref{sec:glo}).
\end{table}

\section{Discussion and Limitations}\label{sec:discussion}

QAD-MultiGuard achieves competitive on-device fraud-detection performance under low-bit quantisation. Nevertheless, several limitations should be considered when interpreting the reported findings and assessing deployment feasibility.

\textbf{Privacy--utility trade-off.}
The 128-dimensional acoustic embedding $\mathbf{F}_v$ demonstrates robustness against reconstruction under the evaluated attack setting ($\text{WER} \ge 0.95$ under white-box GLO attacks, $\text{PESQ} \le 1.21$), which is aligned with common data minimisation objectives in privacy-preserving machine-learning systems. This compression inherently bounds representational capacity, as reflected by the $0.008$ $\text{F}_1$ gap between the BF16 teacher baseline and the full pipeline ($0.931$ versus $0.923$).

An optional $(\epsilon, \delta)$-local differential privacy (LDP) module is supported for stricter privacy requirements and remains disabled in all reported main results. When activated on the edge device using a Gaussian mechanism ($\sigma = 1.0$, corresponding to $\epsilon = 1.5$ at $\delta = 10^{-5}$), performance on TAF-28k decreases from $0.923$ to $0.902$, corresponding to a $0.021$ absolute drop, while introducing approximately 3~\text{ms} of additional latency. Since perturbation is applied once to a 128-dimensional vector, the end-to-end P50 alert latency remains stable at approximately 271~\text{ms} (Figure~\ref{fig4}c). These results characterise a single operating point of the privacy--utility curve; alternative privacy budgets would shift this trade-off accordingly. The reported $\epsilon$ value is an engineering estimate obtained without a certified sensitivity bound (the Gaussian noise is applied to unclipped hidden states), not a full differential-privacy analysis; the LDP module is therefore treated as an optional engineering experiment rather than a formal privacy mechanism.

\textbf{Unlinkability and implicit fingerprint risk.}
The practical evaluation results in Section~\ref{sec:glo} indicate that speech content and speaker identity remain unrecovered from $\mathbf{F}_v$ under the evaluated attack settings. They do not, however, guarantee unlinkability.

Because $\mathbf{F}_v$ is a deterministic function of the acoustic input, repeated observations from the same speaker under similar conditions yield stable embeddings. Consequently, an honest-but-curious cloud (Tier-2) retaining historical representations may cluster embeddings via cosine similarity and link repeated interactions without reconstructing the underlying audio. This differs from reconstruction ($\mathcal{G}_1$) and closed-set speaker identification, as it relies on metric consistency in representation space rather than explicit inference tasks.

The current system does not explicitly optimise for unlinkability, as detection performance and deployment efficiency are prioritised under edge constraints. A systematic quantification of linkability across privacy budgets is not included in this study and represents a necessary extension for deployment-oriented analysis; a promising direction is to inject a session-specific dynamic perturbation or apply an orthogonal rotation matrix to $\bm{F}_v$ that breaks metric consistency across repeated observations while leaving the classification boundary unchanged.

\textbf{Misclassification and robustness.}
Analysis of the $213$ misclassified samples identifies three dominant failure modes: (1) short utterances ($< 8~\text{s}$) account for $41\%$ of errors, limiting the temporal context available for acoustic modelling; (2) white-box PGD-20 perturbations applied to the acoustic input ($\epsilon = 0.10$) cause measurable performance degradation, suggesting sensitivity to adversarial acoustic inputs; and (3) severe background noise ($\text{SNR} < 5\,\text{dB}$) results in a $4.2\%$ drop in discrimination performance. The multimodal fusion architecture partially mitigates these effects by leveraging complementary modalities; however, robustness under adaptive multi-modal attacks remains an open challenge. Likewise, the identity-fraud spoofing threat $\mathcal{G}_3$ is addressed architecturally through cross-modal consistency verification rather than benchmarked against a dedicated voice-spoofing corpus, so its quantitative robustness under realistic spoofing conditions has not been independently verified.

\textbf{Societal consequences of misclassification.}
Beyond aggregate accuracy, the asymmetric consequences of detection errors warrant explicit attention. A false negative allows a fraudulent interaction to proceed, disproportionately harming elderly users and others with lower digital literacy, who are the principal targets of telecommunication fraud. Conversely, a false positive interrupts a legitimate call, imposing a distinct harm---the wrongful freezing of a genuine financial or familial transaction---with particular severity in time-critical situations such as a caregiver's emergency call. The three-tier design mitigates but does not eliminate this tension: the on-device path prioritises recall for early intervention, while the asynchronous cloud review reduces false positives off the critical path. Quantifying these harms, and aligning the decision threshold with stakeholder-defined cost asymmetries, is deferred to a deployment-oriented study conducted with operator and regulator involvement.

\textbf{Deployment-model limitation.} The on-device privacy property assumes that the detection software runs on the data subject's own device; under an operator-side deployment, raw audio is already available in the network and the ``raw audio stays on-device'' property does not apply in the same sense. The present study does not analyse operator-side deployment, data-controller designation, or the corresponding legal basis under PIPL; these are left to deployment-oriented work conducted with operator and regulator involvement.

\textbf{Modality-contribution limitation.} The present evaluation reports only the fused multimodal results on \textit{TAF-28k}; single-modality (text-only and audio-only) baselines are not reported in the tables, so the marginal contribution of the acoustic branch to the final decision is not separately quantified. The learned fusion weights ($w_{\mathrm{text}} = 0.40$, $w_{\mathrm{audio}} = 0.30$; Eq.~\eqref{eq:w-deploy}) suggest that the acoustic branch may be the secondary contributor, which would temper the strength of the ``strong multimodal coupling'' motivation advanced in Section~\ref{sec:related}. A single-modality ablation isolating the acoustic embedding's incremental value has been implemented in the public codebase (text-only / audio-only / fused, \texttt{exp15}), and its measured results are being added in the reproduction run.

\textbf{Primary-corpus limitation and generalisation risk.}
The evaluation is primarily based on the TAF-28k corpus, which was collected under a controlled re-enactment protocol; complementary assessments on AdvFraud-3k and the cross-domain ChiFraud benchmark are also reported. The reliance on a single primary corpus for the main multimodal results nonetheless constitutes a key limitation of the present study, as real-world deployment introduces channel mismatch, regional dialect variation, and evolving fraud strategies that are not fully captured by the current benchmarks.

The $0.011$ $\text{F}_1$ drop observed on ChiFraud (OOD, Section~\ref{sec:crossdata}) provides preliminary evidence of cross-domain sensitivity, although ChiFraud remains a text-only benchmark and does not evaluate the full multimodal pipeline. Several components partially address this limitation, including homologous-teacher distillation. A pilot deployment in a real telecommunication-fraud setting is identified as future work (Table~\ref{tab1}) and has not been completed within the scope of this study. Until field validation becomes available, all reported results should be interpreted as controlled-environment performance. The absence of publicly available multimodal fraud datasets continues to limit cross-study comparability and remains a broader methodological challenge in this research domain.

\textbf{Dataset availability.}
The raw TeleAntiFraud-28k audio used in this study is publicly available via a HuggingFace bucket (\url{https://huggingface.co/buckets/wangdajin062/TeleAntiFraud-bucket}); the original dataset is due to Ma et al.~\citep{1}. The AdvFraud-3k adversarial benchmark is an in-house construction derived from the TAF-28k corpus, with its generation protocol (adversarial rewriting by three annotators under a fourth-reviewer consensus check, Cohen's $\kappa=0.87$) documented in the repository. The QAD training code, configuration files, and experiment logs will be made publicly available upon publication.

Future work will examine training-time privacy-preserving optimisation (e.g., DP-SGD) and compare it with the inference-time local differential privacy mechanism explored as an optional configuration in this study. Taken together, these observations delineate the operational envelope of the proposed framework while identifying key challenges that remain to be addressed.

\section{Conclusion}\label{sec:conclusion}

This study presented QAD-MultiGuard, an edge--cloud collaborative multimodal framework for telecommunication fraud detection under resource-constrained deployment. Motivated by the three coupled constraints of privacy, fidelity, and responsiveness, the framework integrates four components, each addressing a specific constraint: (1)~a pure-KL quantisation-aware distillation strategy that recovers the accuracy lost to 4-bit quantisation (fidelity); (2)~the OV-Freeze regularisation scheme that further stabilises quantisation-sensitive projection layers (fidelity); (3)~a 128-dimensional acoustic embedding with empirical resistance to reconstruction, allowing raw audio to remain on-device (privacy); and (4)~a domain-adapted speculative decoding module that accelerates the asynchronous cloud-side chain-of-thought review ($3.32\times$ wall-clock speedup), complemented by a quantised on-device student sustaining sub-500~ms per-window latency (responsiveness).

Experimental evaluation on TAF-28k indicates that the on-device \texttt{Q4\_K\_M} student retains $98.5\%$ of the unquantised BF16 teacher performance ($\text{F}_1 = 0.917$ versus $0.931$) at a median per-window latency of 268~ms, while the cloud-side \texttt{NVFP4} track reaches $99.1\%$ ($\text{F}_1 = 0.923$) with asynchronous chain-of-thought review, indicating feasibility for real-time edge deployment under the evaluated benchmark conditions. As noted in Section~\ref{sec:experiments}, these tables are drawn from a formal H100 run under a numerical-behaviour-emulation protocol (Section~\ref{sec:nbe}), and a reproduction run aligning the public repository with the reported numbers is in progress. The three constraints are met jointly: the acoustic embedding keeps raw biometric data on-device, the distillation pipeline recovers quantisation-induced degradation, and the on-device quantised student plus lightweight linear fusion sustains sub-500~ms per-window latency.

\textbf{Design takeaways.} Three deployment principles follow directly from these results: (i) keep raw biometric data on-device by emitting only a reconstruction-resistant embedding; (ii) recover low-bit accuracy loss through homologous self-distillation plus variance regularisation rather than post-hoc calibration; and (iii) meet the latency budget on the device with a quantised student and lightweight linear fusion, reserving asynchronous cloud review for off-critical-path refinement.

The primary limitation of this study is its reliance on a single primary corpus (TAF-28k) for the main multimodal evaluation, complemented by adversarial and cross-domain assessments on AdvFraud-3k and ChiFraud. Although these results suggest partial cross-domain robustness, generalisation to unconstrained real-world environments has not yet been formally validated. Short utterances, adversarial perturbations, and distribution shifts remain key challenges for real-world deployment. Addressing these challenges through broader field validation, quantitative privacy--linkability analysis, and continued robustness improvement under realistic operational conditions defines the roadmap for future work.

\section*{CRediT authorship contribution statement}

\textbf{Dajin Wang:} Conceptualization, Methodology, Formal analysis, Writing -- original draft, Writing -- review \& editing.

\textbf{Jianming Bai:} Conceptualization, Methodology, Supervision, Writing -- review \& editing, Funding acquisition.

\textbf{Wu Zhang:} Methodology, Software, Validation, Writing -- review \& editing.

\textbf{Wenbin Guo:} Investigation, Data curation, Validation.

\textbf{Qian Zhao:} Investigation, Data curation, Validation.

\textbf{Liangdong Yang:} Investigation, Data curation, Validation.

\section*{Declaration of generative AI in scientific writing}

During the preparation of this work, the authors utilised generative AI and AI-assisted technologies solely to improve the readability and language of the manuscript. This process was conducted in conjunction with professional linguistic editing. Following the use of these tools, the authors carefully reviewed, verified, and revised the content as necessary, and take full responsibility for the integrity and accuracy of the final published work.

\section*{Declaration of competing interest}

The authors declare that they have no known competing financial interests or personal relationships that could have appeared to influence the work reported in this paper.

\section*{Data availability}

The QAD-MultiGuard training code, configuration files, reproduction scripts, and the anonymised academic text-feature benchmark will be made publicly available upon publication at:

\url{https://github.com/wangdajin062/QAD-MultiGuard}

Researchers interested in reproducing the reported experiments or accessing experimental artefacts prior to formal publication may request early access from the corresponding author.

\section*{Acknowledgments}

This work was supported by the Higher Education Teachers' Innovation Fund Project of the Gansu Provincial Department of Education, China (Grant No.~2023A-224), the National Natural Science Foundation of China (Grant No.~72273059), and the National Key Research and Development Program of China (Grant No.~2024YFB3311600).

The authors sincerely thank Elsevier Language Editing Services for professional English language editing and proofreading support (Order reference: ASLEEX1132254).

\section*{Language editing certificate}

\noindent
\textbf{Certificate of Elsevier Language Editing Services}

\vspace{0.6em}

\renewcommand{\arraystretch}{1.2}

\begin{tabular}{p{4.2cm} p{10cm}}
\textbf{Article Title:} &
\textit{QAD-MultiGuard: An Edge--Cloud Framework for Multimodal Fraud Detection and Risk Assessment} \\
\textbf{Ordered by:} &
Dajin Wang \\
\textbf{Delivery Date:} &
2026-06-19 \\
\textbf{Order Reference:} &
ASLEEX1132254 \\
\end{tabular}

\appendix

\section{Analysis of OV-Freeze and Acoustic Embedding Reconstruction Resistance}
\label{app:ovf}

%------------------------------------------------
\subsection{Boundedness of the OV-Freeze Correction Factor}
\label{app:bounded}

The boundedness of the OV-Freeze correction factor $c_\ell$ is analysed to characterise its effect on the numerical stability of intermediate activation rescaling and to clarify the underlying conditions used in the subsequent optimisation analysis.

\medskip
\noindent\textbf{Lemma A.1 (Boundedness of the correction factor).}
Let $\mathbf{y}_\ell \in \mathbb{R}^{N \times d}$ denote the intermediate activation at layer $\ell$, and define the rescaled activation as
\[
\mathbf{y}'_\ell = c_\ell \mathbf{y}_\ell, \quad
c_\ell = \sqrt{\frac{\sigma^2_{\mathrm{BF16},\ell}}{\mathrm{Var}(\mathbf{y}_\ell)+\epsilon}}.
\]
Then there exist constants $0 < c_{\min} \le c_{\max} < \infty$ such that
\[
c_\ell \in [c_{\min}, c_{\max}]
\]
whenever the activation variance $\mathrm{Var}(\mathbf{y}_\ell)$ is finite.

\medskip
\noindent\textbf{Proof.}
Since $\sigma^2_{\mathrm{BF16},\ell} > 0$ is fixed, the activation variance $\mathrm{Var}(\mathbf{y}_\ell)$ is finite, and $\epsilon > 0$, the denominator remains strictly positive. Therefore, $0 < c_\ell < \infty$ holds. Over a finite training trajectory with bounded activation variance, the infimum and supremum
\[
c_{\min} = \inf_{\ell} c_\ell, \qquad c_{\max} = \sup_{\ell} c_\ell,
\]
are attained and finite, and directly satisfy the inequality
\[
0 < c_{\min} \le c_\ell \le c_{\max} < \infty.
\]
This establishes the claimed boundedness property and completes the proof.

\medskip
\noindent\textbf{Remark 1.}
Under the detached-gradient formulation specified in the joint objective, the correction factor $c_\ell$ operates as a constant scaling matrix during backward error propagation. This mechanism mitigates gradient explosion while preserving structural directionality within established bounded limits, thereby supporting numerical stability without altering the underlying optimisation dynamics.

%------------------------------------------------
\subsection{Optimisation Stability of OV-Freeze}
\label{app:lipschitz}

We analyse the effect of OV-Freeze on the backward gradient flow, showing that the stop-gradient rescaling keeps the effective gradient bounded and does not introduce unbounded amplification.

\medskip
\noindent\textbf{Proposition A.2 (Bounded gradient rescaling under OV-Freeze).}
Under the stop-gradient formulation of Eq.~\eqref{eq:ovf-detach-en}, the backward gradient at layer $\ell$ is scaled by the bounded factor $c_\ell \in [c_{\min}, c_{\max}]$ established in Lemma~A.1:
\[
\frac{\partial L_{\mathrm{QAD}}}{\partial\bm{y}_\ell}
= c_\ell \cdot \frac{\partial L_{\mathrm{QAD}}}{\partial\bm{y}'_\ell},
\]
and therefore the gradient norm is amplified by at most a constant factor:
\[
\left\|\frac{\partial L_{\mathrm{QAD}}}{\partial\bm{y}_\ell}\right\|
\;\le\; c_{\max}\,\left\|\frac{\partial L_{\mathrm{QAD}}}{\partial\bm{y}'_\ell}\right\|.
\]
Because $c_\ell$ is detached from the computational graph, it acts as a constant in backward propagation, so the rescaling mitigates the risk of gradient explosion without altering the direction of the optimisation update.

\medskip
\noindent\textbf{Remark 2.}
The variance-matching loss $L_{\mathrm{OVF}}$ of Eq.~\eqref{eq:ovf-loss} separately steers the parameters toward low-variance regions through its own gradient. The analysis above establishes only the boundedness of the per-layer gradient rescaling; we do not claim a formal convergence-rate guarantee for the joint objective, as $c_\ell$ depends on the network outputs and hence on the parameters.

%------------------------------------------------
\subsection{Empirical Reconstruction Resistance of $\mathbf{F}_v$}
\label{app:nonrecon}

The empirical resistance of the $128$-dimensional acoustic embedding $\mathbf{F}_v$ against reconstruction is evaluated under white-box GLO and black-box inversion attacks. Across all evaluated settings, the reconstructed speech signals yield the following empirical bounds:
\[
\text{WER} \ge 0.95, \quad \text{PESQ} \le 1.21, \quad
\text{Speaker-ID accuracy} \le 8.3\%.
\]

These metrics indicate that $\mathbf{F}_v$ resists intelligible reconstruction under the evaluated attack models, providing empirical evidence rather than a formal information-theoretic guarantee.

\medskip
\noindent\textbf{Remark 3.}
The empirical reconstruction resistance underscores that while the deterministic encoder guarantees a non-zero mutual information condition $I(x;\mathbf{F}_v) > 0$, the remaining representation entropy is insufficient to resolve sub-frame acoustic trajectories under the evaluated adversarial inversion pipelines. These findings define the empirical operating boundary of the proposed representation under the evaluated threat models and do not exclude stronger adversarial strategies outside the considered evaluation scope.


\bibliographystyle{cas-model2-names}
\bibliography{ref_v4}

\end{document}
